\documentclass[aic]{iosart2x}

%----Generally useful
\usepackage{array}
\usepackage{hyperref}
\usepackage{wasysym}
\usepackage{amsfonts}
\usepackage{amsmath}
% \usepackage{graphicx}
% \usepackage{algorithm2e}
\newcommand{\smalltt}[1]{\small \texttt{#1}}
\newcommand{\tildenot}{\raisebox{0.4ex}{\texttildelow}}

%----Useful for this document
\usepackage{dcolumn}
\usepackage{rotating}
\usepackage{pdflscape}
\newcommand{\bb}[1]{\bf {#1}}
\newcommand{\xxx}{$\times$}
\newcommand{\PreserveBackslash}[1]{\let\temp=\\#1\let\\=\temp}
\newcolumntype{C}[1]{>{\PreserveBackslash\centering}p{#1}}
\newcolumntype{R}[1]{>{\PreserveBackslash\raggedleft}p{#1}}
\newcolumntype{L}[1]{>{\PreserveBackslash\raggedright}p{#1}}

%----Making things more compact
\newenvironment{packed_enumerate}{
\vspace*{-0.3em}
\begin{enumerate}
  \setlength{\partopsep}{0pt}
  \setlength{\itemsep}{1pt}
  \setlength{\parskip}{0pt}
  \setlength{\parsep}{0pt}
}{\end{enumerate}}
\newenvironment{packed_itemize}{
\vspace*{-0.3em}
\begin{itemize}
  \setlength{\partopsep}{0pt}
  \setlength{\itemsep}{1pt}
  \setlength{\parskip}{0pt}
  \setlength{\parsep}{0pt}
}{\end{itemize}}
\newcommand{\ass}[1]{\footnotesize {#1}}
%----Suppress extra space in texttt mode
\AddToHook{cmd/ttfamily/after}{\frenchspacing}
% \renewcommand{\arraystretch}{0.8}
\renewcommand{\floatpagefraction}{0.9}
% \renewcommand{\textfloatsep}{2.0ex}
% \setlength{\floatsep}{2.0pt plus 2.0pt minus 2.0pt}
% \setlength{\textfloatsep}{5.0pt plus 2.0pt minus 0.0pt}
% \renewcommand{\dbltextfloatsep}{2.0ex}
% \renewcommand{\textfraction}{0.07}
% \renewcommand{\topfraction}{0.9}
% \renewcommand{\bottomfraction}{0.9}
% \setlength{\tabcolsep}{3pt}
%% One can fix some overfulls
\sloppy

\newcommand{\LAST}{\mbox{CASC-30}}
\newcommand{\THIS}{\mbox{CASC-J13}}
\newcommand{\NEXT}{\mbox{CASC-31}}
\newcommand{\NUMBEROFENTRANTS}{\mbox{Twenty three}}
\newcommand{\NTHCOMPETITION}{\mbox{thirty-first}}

\begin{document}
\begin{frontmatter}                           % The preamble begins here.
%
%\pretitle{Pretitle}
\title{The 13th IJCAR Automated Theorem Proving System Competition -- \THIS{}}
\runtitle{\THIS{}}
%\subtitle{Subtitle}
\maketitle
%
\begin{aug}
\author[A]{\fnms{Geoff} \snm{Sutcliffe} \ead[label=e1]{geoff@tptp.org} \thanks{Corresponding author. \printead{e1}.}}
\address[A]{Department of Computer Science, \orgname{University of Miami},
\cny{USA}\printead[presep={\\}]{e1}}
\end{aug}
\runauthor{Sutcliffe}

\begin{abstract}
The CADE ATP System Competition (CASC) is the annual evaluation of fully automatic, classical
logic, Automated Theorem Proving (ATP) systems –- the world championship for such systems.
\THIS{} was the \NTHCOMPETITION{} competition in the CASC series.
\NUMBEROFENTRANTS{} ATP systems competed in the various divisions.
This paper presents an outline of the competition design and a commentated summary of the results.
\end{abstract}

\begin{keyword}
automated theorem proving\sep competition
\end{keyword}
%
\end{frontmatter}
%--------------------------------------------------------------------------------------------------
\paragraph{Author contributions:} The paper, except for Section~\ref{Systems}, was written by 
the paper author.
The three subsections of Section~\ref{Systems} were written by Nicolas Murzin (FindProof), 
Olivier Roland (mrs), and Theodore Kim (SUPr).

\paragraph{Statements and Declarations:} 
Ethical considerations - Not applicable.
Consent to participate - Not applicable.
Consent for publication - Not applicable.
Conflicting interest - The author(s) declare no potential conflicts of interest with respect to 
the research, authorship, and/or publication of this article.
Funding statement - Not applicable.
%--------------------------------------------------------------------------------------------------
\section{Introduction}
\label{Introduction}

Automated Theorem Proving (ATP) deals with the task of proving theorems from axioms – the
derivation of conclusions that follow inevitably from known facts~\cite{RV01-HAR}.
The converse task of disproving conjectures is another facet of interest~\cite{CS03,BN10-ITP}.
The axioms and conjecture are written in an appropriately expressive logic, and the solutions
(proofs and models) are often similarly written in logic~\cite{Sut23-IGPL}.
The CADE ATP System Competition (CASC)~\cite{Sut16} is the annual evaluation of fully automatic,
classical logic, ATP systems – the world championship for such systems.
One purpose of CASC is to provide a public evaluation of the relative capabilities of ATP systems.
CASC evaluates the performance of the ATP systems in terms of
the number of problems solved,
the number of acceptable solutions (proofs and models) output, and
the average time taken for problems solved,
in the context of
a bounded number of eligible problems and
specified time limits.
Additionally, CASC aims to
stimulate ATP research,
motivate development and implementation of robust ATP systems that can be easily and usefully
deployed in applications,
provide an inspiring environment for personal interaction between ATP researchers,
and
expose ATP systems within and beyond the ATP community.

CASC is held at each CADE (the International Conference on Automated Deduction) and IJCAR
(the International Joint Conference on Automated Reasoning) conference -- the major forums
for the presentation of new research in all aspects of automated deduction.
\THIS{} was held on 27th July 2026, as part of the
% 30th International Conference on Automated Deduction (CADE-30).
13th International Joint Conference on Automated Reasoning (IJCAR 2024)\footnote{%
CADE was a constituent conference of the 13th IJCAR, hence CASC-``J13'',
which in turn was part of the Federated Logic Conference 2026},
in Lisbon, Portugal.
It was the \NTHCOMPETITION{} competition in the CASC series; see~\cite{Sut26-CASC} and citations
therein, and the CASC web site \href{http://tptp.org/CASC}{\smalltt{tptp.org/CASC}}, for 
information about previous CASCs.
The ``FLoC Olympic Games'' hosted fifteen competitions in various areas of logic and reasoning, 
which provided the \THIS{} entrants with opportunities to observe and interact with developers 
of other types of reasoning tools.
% , and~\cite{SS06-SoCASC} for a historical overview and analysis of the
% first ten CASCs.

\THIS{} was organized by Geoff Sutcliffe,
% assisted by Martin Desharnais for the SLH division,
and overseen by a panel consisting of
Aart Middeldorp, Cl{\' a}udia Nalon, and Martina Seidl.
The CASC panel has a few, but important, responsibilities:
\begin{packed_itemize}
\item Advise (by email) on contentious issues that arise before the competition, and that cannot
      be agreeably resolved by the organizer and entrants.
\item With the help of the organizer, adjudicate on the acceptability of the sample proofs and
      models that are submitted before the competition.
\item Supply a seed digit for the random problem selection process.
\item If possible, be physically present at the start of the competition, so entrants can see who
      they are and raise any issues they are concerned about.
\item Be contactable during the competition to advise on contentious issues that arise and that
      cannot be agreeably resolved by the organizer and entrants.
\item With the help of the organizer, adjudicate on the acceptability of the winners' proofs and
      models.
\item If possible, be physically present at the end of the competition for resolving possible 
      open issues, confirming the final system ranking in each division, and assisting in winner 
      announcements.
\end{packed_itemize}

\THIS{} was run on computers provided by the StarExec project~\cite{SST14} at the
University of Miami.
% the Max-Planck-Institut f{\"u}r Informatik, Saar\-br{\"u}cken, Germany.
% and by the Department of Computer Science at the University of Manchester.
% the Department of Computer Science at the University of Miami.
The \THIS{} web site \href{http://tptp.org/CASC/J13}{\smalltt{tptp.org/CASC/J13}}
provides access to all the resources used before, during, and after the event.

The design and organization of CASC has evolved over the years to a sophisticated state.
An outline of the \THIS{} design and organization is provided here.
The details are in~\cite{Sut26-CASC-J13} and on the \THIS{} web site.
Important changes for \THIS{} were
(for readers already familiar with the general design of CASC):
\begin{packed_itemize}
\item The TFA and EPR divisions went on hiatus.
\item The FNT division returned from hiatus.
\item Inference steps in proofs were expected to list only parents that are not used in the 
      inference.
\item Proofs were expected to graft in subproofs from external systems.
\item A new performance measure, proof granularity, was recorded and reported.
\item The ProoVer competition for proof verifiers was held for the first time in 2026, 
      as a companion competition to \THIS{}.
\end{packed_itemize}
The CASC rules, specifications, and deadlines are absolute.
Only the panel has the right to make exceptions.
It is assumed that all entrants have read the documentation related to the competition, and have
complied with the competition rules.
Non-compliance with the rules can lead to disqualification.
A ``catch-all'' rule is used to deal with any unforeseen circumstances: 
{\em No cheating is allowed}.
The panel is allowed to disqualify entrants due to unfairness, and to adjust the competition
rules in case of misuse.

\vspace*{0.5em}
The rest of this paper is organized as follows:
Section~\ref{DivisionsSystems} describes the competition divisions and the ATP systems that
entered the various divisions.
Sections~\ref{Infrastructure} and \ref{Registration} describe the competition infrastructure
and the requirements for the ATP systems.
Section~\ref{Evaluation} describes how the systems are evaluated.
Section~\ref{Results} provides a commentated summary of the results.
Section~\ref{Systems} contains short descriptions of three of the ATP systems.
% the winning systems.
Section~\ref{Acceptable} discusses the increasingly precise definition and requirements for
acceptable proofs in TPTP format, in preparation for \NEXT{}.
Section~\ref{ProoVer} gives an overview of the companion ProoVer competition, and its ongoing
relationship with CASC.
Section~\ref{Conclusion} concludes and discusses plans for future CASCs.

\vspace*{0.5em}
\emph{Executive summary:} \THIS{} was the \NTHCOMPETITION{} annual evaluation of fully 
automatic, classical logic, Automated Theorem Proving (ATP) systems.
\THIS{} fulfilled its objectives by evaluating the relative capabilities of ATP systems,
and stimulating development and interest in ATP.
The highlights of \THIS{} were many new entrants, and the companion ProoVer competition.
Despite the continued dominance of the Vampire systems, the stimulating and interactive nature
of CASC has kept the developer and user communities engaged.
The experience gained from \THIS{} has prompted further focus on proof production and verification
in \NEXT{}.

\vspace*{0.5em}
\emph{A Tense Note:}
Attentive readers will notice changes between the present and past tenses in this paper.
Many parts of CASC are established and stable -- they are described in the present tense (the
rules are the rules).
Aspects that were particular to \THIS{} are described in the past tense so that they make sense
when reading this after the event.

%--------------------------------------------------------------------------------------------------
\section{Divisions and Systems}
\label{DivisionsSystems}

CASC is divided into divisions according to problem and system characteristics, in a coarse
version of the TPTP problem library's Specialist Problem Classes (SPCs)~\cite{SS01}.
Each division uses problems that have certain logical, language, and syntactic characteristics,
so that the ATP systems that compete in a division are, in principle, able to attempt all the 
problems in the division.
Some divisions are further divided into problem categories that make it possible to analyze, at
a more fine-grained level, which systems work well for what types of problems.
Table~\ref{DivisionsTable} catalogs the divisions and problem categories of \THIS{}.
The example problems can be viewed online at
\href{https://tptp.org/cgi-bin/SeeTPTP?Category=Problems}{\smalltt{tptp.org/cgi-bin/SeeTPTP?Category=Problems}}.
Section~\ref{TPTPProblems}
% , \ref{SLHProblems}, and~\ref{ICUProblems} 
explains what problems are eligible for use in each division and category.
The competition rankings are at (only) the division level, as explained in 
Section~\ref{Evaluation}.

\begin{table*}[htbp]
\caption{Divisions and Problem categories}
\label{DivisionsTable}
\begin{tabular}{p{1cm}p{7cm}p{5cm}}
\hline
Division  & Problems              & Problem categories \\
\hline
\textbf{THF} &
\textbf{T}yped (monomorphic) \textbf{H}igher-order \textbf{F}orm theorems
(axioms with a provable conjecture). &
\textbf{TNE} -- \textbf{T}HF with \textbf{N}o \textbf{E}quality

      Example: \texttt{NUM738\string^{}1}

\textbf{TEQ} -- \textbf{T}HF with \textbf{EQ}uality

      Example: \texttt{SET171\string^{}3} \\
\hline
% \textbf{THN} &
% \textbf{T}yped \textbf{H}igher-order form \textbf{N}on-theorems
% (axioms with a countersatisfiable conjecture, and satisfiable axiom sets without a conjecture). &
% \textbf{TNN} -- \textbf{T}H\textbf{N} with \textbf{N}o equality
%
%       Example: \texttt{CSR142\string^{}3}
%
% \textbf{TNQ} -- \textbf{T}H\textbf{N} with e\textbf{Q}uality
%
%       Example: \texttt{PHI003\string^{}3} \\
% \hline
% \textbf{TFA} &
% \textbf{T}yped (monomorphic) \textbf{F}irst-order form theorems
%       with \textbf{A}rithmetic (axioms with a provable conjecture). &
% \textbf{TFI}~ -- \textbf{TF}A with only \textbf{I}nteger arithmetic
%
%       Example: \texttt{DAT016\_1}
%
% % \textbf{TFR} -- \textbf{TF}A with only \textbf{R}ational arithmetic
% 
% %       Example: \texttt{ARI621\_2}
% 
% \textbf{TFE} -- \textbf{TF}A with only r\textbf{E}al arithmetic,
%       e.g., \texttt{MSC022\_2} \\
% \hline
% \textbf{TFN} & \textbf{T}yped \textbf{F}irst-order form \textbf{N}on-theorems
%     (axioms with a countersatisfiable conjecture, and satisfiable axiom sets without a conjecture)
%     without arithmetic. &
%     Example: \texttt{COM002\_20} \\
% \hline
\textbf{FOF} &
\textbf{F}irst-\textbf{O}rder \textbf{F}orm theorems (axioms with a provable
    conjecture). &
\textbf{FNE} -- \textbf{F}OF with \textbf{N}o \textbf{E}quality.

    Example: \texttt{COM003+1}

\textbf{FEQ} -- \textbf{F}OF with \textbf{EQ}uality.

    Example: \texttt{SEU147+3} \\
\hline
\textbf{FNT} &
\textbf{F}OF \textbf{N}on-\textbf{T}heorems (axioms with a countersatisfiable
    conjecture, and satisfiable axioms without a conjecture). &
\textbf{FNN} -- \textbf{FN}T with \textbf{N}o equality.

    Example: \texttt{KRS173+1}

\textbf{FNQ} -- \textbf{FN}T with e\textbf{Q}uality.

    Example: \texttt{MGT033+2} \\
\hline
% \textbf{EPR} &
% \textbf{E}ffectively \textbf{PR}opositional clause normal form theorems and non-theorems (clause
%     sets). &
% \textbf{EPU} -- \textbf{EP}R \textbf{U}nsatisfiable clause sets.
% 
%     Example: \texttt{GEO079-1}
% 
% \textbf{EPS} -- \textbf{EP}R \textbf{S}atisfiable clause sets.
%
%     Example: \texttt{PUZ001-3} \\
% \hline
\textbf{UEQ} &
     \textbf{U}nit \textbf{EQ}uality theorems in clause normal form
     (unsatisfiable clause sets). &
     Example: \texttt{RNG026-7} \\
\hline
% \textbf{SLH} &
% Typed (monomorphic) higher-order theorems (axioms with a provable conjecture),
% generated by Isabelle's \textbf{SL}edge\textbf{H}ammer system~\cite{PB10}.  &
% The problems were generated from sessions in Isabelle's Archive of Formal Proofs~\cite{BH+15}. \\
% \hline
% \textbf{ICU} &
% First-order theorems (axioms with a provable conjecture) provided by the entrants. &
% The problems supplied by an entrant were expected to be easy enough for that entrant's ATP system,
% but difficult for the other entrants' systems, i.e., each entrant is saying to the others: ``I
% Challenge yoU!''. \\
% \hline
% \textbf{LTB} &
% Theorems (axioms with a provable conjecture) from \textbf{L}arge \textbf{T}heories, presented
% in \textbf{B}atches.
% A large theory typically has many functions and predicates, and many axioms of which only a few
% are required for a proof of a theorem.
% The problems in a batch are given to an ATP system all at once, and typically have a common core
% set of axioms.
% The batch presentation allows the ATP systems to load and preprocess the common core set of
% axioms just once, and to share logical and control results between proof searches.
% Each problem category might be accompanied by a set of training problems and their solutions,
% % taken from the same source as the competition problems,
% which can be used for ATP system tuning during (typically at the start of) the competition &
% \textbf{VBT} -- Problems generated by Isabelle's Sledgehammer system, from the
% ``{\bf v}an Emde {\bf B}oas {\bf T}rees'' entry in the Archive of Formal Proofs.
% See Section~\ref{LTBProblems} for details.
% There was no common core set of axioms.
%
% \vspace{0.7em}
% Six versions of each VBT problem were provided:
% % a first-order form (FOF) version,
% monomorphic and polymorphic typed first-order versions,
% monomorphic and polymorphic typed extended first-order versions,
% and
% monomorphic and polymorphic typed higher-order versions without \smalltt{\$ite}, 
% \smalltt{\$let}, or tuples.
% Systems could attempt as many of the versions as they wanted, and could attempt the problems
% and versions in any order including in parallel.
% A solution to any version counted as a solution to the problem.
%
% \vspace{0.7em}
% The VBT category was accompanied by training data. \\
% \hline
\end{tabular}
\end{table*}

ATP systems that cannot run in the competition divisions for any reason, e.g., the system cannot
run on the competition computers, or the entrant is an organizer, can be entered into the 
demonstration division.
The demonstration division uses the same problems as the competition divisions, and the entry
specifies which competition divisions' problems are to be used.
Demonstration division systems can run on the competition computers, or on computers supplied by 
the entrant.
The demonstration division results are presented along with the competition divisions' results,
but might not be comparable with those results.

Table~\ref{Entrants} lists the 23 
% \NUMBEROFENTRANTS{} 
ATP systems that competed in \THIS{}, which divisions they entered, and their entrants' 
information.
A division acronym in {\em italics} indicates the system was in the demonstration division.
The division winners of the previous CASC (\LAST{}) and the Prover9~1109a system are automatically
entered into the demonstration division to provide benchmarks against which progress can be judged.
System descriptions are in the competition proceedings~\cite{Sut24-CASC-J12} and on the \THIS{}
web site.

\begin{table*}
\setlength{\tabcolsep}{5pt}
\footnotesize
\caption{The ATP systems and entrants}
\label{Entrants}
\begin{tabular}{lp{1.5em}p{1.5em}p{1.5em}p{1.5em}>{\raggedright\arraybackslash}p{4.4cm}>{\raggedright\arraybackslash}p{3.4cm}}
ATP System                & \multicolumn{4}{l}{Divisions} & Entrant (Associates)              & Entrant's Affiliation\\
\hline
Connect++ 0.7.2           &   &FOF&FNT&   & Dr Sean B Holden                                  & University of Cambridge \\
CSI++ 1.0                 &   &FOF&   &UEQ& Guoyan Zeng (Yang Xu, Baojie Rui, Stephan Schulz, Peiyao Liu, Hailin Guo, Jun Liu, Shuwei Chen)
                                                                                              & Xihua University \\
Drodi 4.1.1               &   &FOF&FNT&UEQ& Oscar Contreras                                   & Amateur Programmer \\
E 3.5.1                   &THF&FOF&FNT&UEQ& Stephan Schulz                                    & DHBW Stuttgart \\
FindProof 0.1             &   &FOF&   &UEQ& Nik Murzin                                        & Wolfram Institute \\
FMB4J 0.1                 &   &   &FNT&   & Michael Rawson (Joe Hall)                         & University of Southampton \\
iProver 3.9.4             &   &FOF&FNT&UEQ& Konstantin Korovin                                & University of Manchester \\
LEO-II 1.7.0              &THF&   &   &   & Alexander Steen (Christoph Benzmüller)            & University of Greifswald \\
Leo-III 1.9.0             &THF&   &   &   & Alexander Steen (Christoph Benzmüller)            & University of Greifswald \\
LisaST 0.9                &   &{\em FOF}&   &   & Simon Guilloud                              & EPFL \\
Mace4 2026-6A             &   &   &FNT&   & Jeff Machado (Larry Lesyna)                       & Independent Researcher \\
mrs 0.2.0                 &   &FOF&   &UEQ& Olivier Roland                                    & Independent Researcher \\
Prover9 1109a             &   &{\em FOF}&   &   & CASC (Bob Veroff, Bill McCune)              & CASC fixed point \\
Prover9 2026-6A           &   &FOF&   &UEQ& Jeff Machado (Larry Lesyna)                       & Independent Researcher \\
SATResetCoP 1.0           &   &FOF&   &   & Martin Fixman (Fredrik Rømming, Sean Holden)      & University of Cambridge \\
SPASS-SCL 0.1.1           &   &FOF&   &   & Simon Schwarz (Yasmine Briefs, Lorenz Leutgeb, Christoph Weidenbach, Fabian Wobito)
                                                                                              & Max Planck Institute for Informatics \\
SUPr 1.0                  &   &FOF&   &   & Teddy Kim (Adam Pease)                            & Naval Postgraduate School \\
Twee 2.7                  &   &   &   &UEQ& Nick Smallbone (Guy Axelrod)                      & Chalmers University of Technology \\
Vampire 4.8               &   &   &{\em FNT}&   & CASC                                        & CASC-29 winner \\
Vampire 5.0               &{\em THF}&{\em FOF}&   &{\em UEQ}& CASC                            & \LAST{} winner \\
Vampire 5.0.1             &THF&FOF&FNT&UEQ& Márton Hajdu (Filip Bártek, Ahmed Bhayat, Jonas Bodingbauer, Robin Coutelier, 
                                                Matthias Hetzenberger, Petra Hozzová, Laura Kovács, Jack McKeown, Merisa Mustajbasic, 
                                                Axel Polaczek, Jakob Rath, Michael Rawson, Giles Reger, Martin Riener, Martin Suda, 
                                                Johannes Schoisswohl, Andrei Voronkov, 
                                                Eva Maria Wagner)                                 & TU Wien \\
VIP 1.71828459        &   &FOF&   &   & Ilies Nokrani (David Delahaye, Clémence Gerardin) & Université Montpellier - LIRMM \\
Zipperp'n 2.1.9999        &THF&FOF&   &   & Jasmin Blanchette
                          (Alexander Bentkamp, Simon Cruanes, Petar Vukmirovi{\'c})               & Ludwig-Maximilians-Universität München\\
\hline
\end{tabular}
\end{table*}

%--------------------------------------------------------------------------------------------------
%----HERE IS WHERE WE START SLIDING INTO PAST TENSE
%--------------------------------------------------------------------------------------------------
\section{Infrastructure}
\label{Infrastructure}

%--------------------------------------------------------------------------------------------------
\subsection{Computers}
\label{Computers}

The StarExec Miami computers used for the competition have
two octa-core Intel(R) Xeon(R) CPU E5-2620~v4 @~2.10GHz CPUs, with hyper-threading (two threads 
per core),
256GiB memory,
the Ubuntu 24.04.3 LTS operating system, with Linux kernel 6.8.0-71-generic.
There were 30 computers available for \THIS{}, i.e., 60 CPUs.
One ATP system runs on one CPU at a time.
StarExec uses Linux's \texttt{sched\_setaffinity} to restrict each system run to a single CPU,
and \texttt{setrlimit} to limit memory use to 128 GiB.
This separation avoids contention that might affect performance measurements~\cite{BLW19,FG+24}.
Systems can use all the cores on the CPU.
In \THIS{} all the demonstration division systems used the competition computers.
The StarExec Miami computers used for CASC are the same as are publicly available to the TPTP 
community, which allows system developers to test and tune their systems in exactly the same 
environment as is used for the competition.

In \THIS{} a total of 1618 CPU hours were used in 453 wall clock hours, i.e., 5.9 hours spread
over the 60 CPUs, which allowed the competition to run live in one conference day.
For the problems that were solved, 185 CPU hours were used in 34 wall clock hours.

%--------------------------------------------------------------------------------------------------
\subsection{Problems for the TPTP-based Divisions}
\label{TPTPProblems}

The problems for CASC
% for the THF, FOF, EPR, and UEQ divisions 
are taken from the latest release of the
Thousands of Problems for Theorem Provers (TPTP) problem library~\cite{Sut17}.
For \THIS{} that was release v9.3.0.
The TPTP version used for CASC is released only after the competition has started, so that new
problems in the release have not been seen by the entrants.
The problems have to meet certain criteria to be eligible for use:
\begin{packed_itemize}
\item The TPTP tags problems that are designed specifically to be suited or ill-suited to some ATP
      system, calculus, or control strategy as \textit{biased}.
      They are excluded from the competition.
\item The problems must be syntactically non-prop\-osi\-tional.
\item The TPTP uses system performance data in the Thousands of Solutions from Theorem Provers
      (TSTP) solution library~\cite{Sut07-CSR} to compute problem difficulty ratings in the 
      range 0.00 (easy) to 1.00 (unsolved)~\cite{SS01}.
      Problems with ratings in the range 0.21 to 0.99 are eligible.
      The upper bound of 0.99 excludes problems that probably would not be solved by any of the 
      systems, and thus would not differentiate between the systems.
      The lower bound of 0.21 was chosen (many years ago, and it has worked successfully) to
      exclude problems that would be solved by most of the systems and thus not differentiate
      between the systems.
      Problems of lesser and greater ratings are made eligible if there are not enough problems
      with ratings in that range -- this was unnecessary in \THIS{}.

      \hspace*{1em}
      Systems can be submitted before the competition so that their performance data is used in
      computing the problem ratings -- problems that are newly solved get a rating less than
      1.00 and thus become eligible (until the rating drops below 0.21).
      The rating calculation also uses performance data from ATP systems that are not
      entered into the competition, which can produce ratings that make some problems eligible
      for selection but easy or unsolvable for the systems in the competition.
      Using problems that are solved by all or none of the competition systems does not
      affect the competition rankings, has the benefit of placing the systems' performances
      in the context of the state-of-the-art in ATP, but does reduce the differentiation between
      the systems in the competition.
\end{packed_itemize}

In order to ensure that no system receives an advantage or disadvantage due to the specific
presentation of the problems in the TPTP problem library, the problems are obfuscated by
% \item rename all predicate and function symbols to meaningless symbols
stripping out all comment lines, including the problem header,
randomly reordering the formulae (\texttt{include} directives are left before the formulae, 
type declarations and definitions are left before the symbols' uses),
randomly swapping the arguments of associative connectives,
randomly reversing implications, and
randomly reversing equalities.

The numbers of problems used in each division and problem category are constrained by the
numbers of eligible problems, the numbers of systems entered in the divisions, the number of 
CPUs available, the time limits, and the time available for running the competition live
in one conference day, i.e., in about 6 hours.
The numbers of problems used are set within these constraints according to the judgement of the 
competition organizer.
The problems used are randomly selected from the eligible problems based on a seed supplied by
the competition panel:
\begin{packed_itemize}
\item The selection is constrained so that no division or category contains an excessive number
      of very similar problems, according to the ``very similar problems'' (VSP) lists distributed
      with the TPTP problem library~\cite{Sut00-JAR}.
      For each problem category in each division, if the category is going to use $N$ problems
      and there are $L$ VSP lists that have an intersection of at least $N/(L+1)$ with the
      eligible problems for the category, then maximally $N/(L+1)$ problems are taken from each
      VSP list.
\item In order to combat excessive tuning towards problems that are in the preceding TPTP version, 
      the selection is biased to select problems that are new in the TPTP version used until 
      50\% of the problems in each problem category have been selected or there are no more new 
      problems to select, after which random selection from old and new problems continues.
      The number of new problems used depends on how many new problems are eligible and the
      limitation on very similar problems.
\item Problems with rating 0.21 to 0.99 are selected before problems with other ratings.
\end{packed_itemize}
Table~\ref{ProblemNumbers} gives the numbers of eligible problems, the maximal numbers that could
be used after taking into account the limitation on very similar problems,
% (in the context of the numbers of problems used),
and the numbers of problems used in each division and category.
As can be seen, the target of 50\% new problems was not met, but there were enough new problems 
in the TEQ, FEQ, and UEQ categories to expose any excessive tuning to the previous TPTP release.
% See Section~\ref{Conclusion} for a plea to the community to submit new problems to the TPTP 
% problem library.

\begin{table*}[tbh]
\caption{Numbers of eligible and used problems}
\label{ProblemNumbers}
\begin{tabular}{l|rr|rr|rr|r}
\hline
\multicolumn{1}{l}{Division} &
\multicolumn{2}{c}{THF} &
\multicolumn{2}{c}{FOF} &
\multicolumn{2}{c}{FNT} &
\multicolumn{1}{c}{UEQ} \\
Category                  & {\tiny TNE} & {\tiny TEQ} 
                                      & {\tiny FNE} & {\tiny FEQ} 
                                                  & {\tiny FNN} & {\tiny FNQ} & \\
\hline
Eligible                  & 147 &1155 & 249 &4135 &  91 & 117 & 439 \\
Usable                    & 139 & 439 & 117 & 912 &  64 & 117 & 439 \\
Used                      & 100 & 300 & 100 & 300 &  50 & 100 & 400 \\
New                       &   0 &  31 &   0 &  79 &   0 &   0 &  21 \\
% \hspace*{2ex}Max new      &
% \hspace*{2ex}Usable       &
% \hspace*{2ex}Usable new   &
\hline
\end{tabular}
\end{table*}

The problems are given to the ATP systems in TPTP format, with {\tt include} directives,
in increasing order of TPTP difficulty rating.

%--------------------------------------------------------------------------------------------------
% \subsection{Problems for the SLH Division}
% \label{SLHProblems}
% 
% For the SLH division of CASC-29, Isabelle's Sledgehammer system~\cite{PB10} was used to generate
% 8400 problems that could be used, of which 1000 appropriately difficult problems were selected 
% based on performance data~\cite{Sut23-CASC-29,SD24-CASC}.
% For the SLH division of \LAST{}, the same problem set was used, and 1000 problems not used in 
% CASC-29 were selected.
% For the SLH division of \THIS{}, the same problem set was used, and 1000 problems not used in
% CASC-29 or \LAST{} were selected.
% The reuse of the problem set was announced in advance of \THIS{}, so that developers could tune
% their systems using the CASC-29 and \LAST{} problems.
% The problems are not modified by any preprocessing, thus allowing the ATP systems to take
% advantage of natural structure that occurs in the problems.
% 
% Of the 1000 problems selected for \THIS{}, 383 had not been solved in the testing done before
% CASC-29.
% If many of them were solved in the competition, that would indicate progress.
% See Section~\ref{SLH} for the results on these problems.
% 
% The problems are given in a roughly estimated increasing order of difficulty.
% 
%--------------------------------------------------------------------------------------------------
% \subsection{Problems for the ICU Division}
% \label{ICUProblems}
% 
% For the ICU division each entrant had to submit 20 FOF theorems (axioms with a provable conjecture).
% The problems had to be provided in decreasing order of desired use in the division, e.g., from 
% hardest to easiest for other systems.
% The problems had to be all different, as assessed by the competition organizer.
% It was expected that each entrant would submit problems that are easy enough for that entrant's
% system, but difficult for the other entrants' systems, i.e., each entrant is saying to the
% others: ``I Challenge yoU!''.
% 
% For \THIS{} the 10 problems chosen for the \LAST{} winner, Vampire~4.9, were reused.
% Then the top 13 problems were taken from each of the entrant's submissions, excluding duplicates.
% That gave a total of 101 problems.
% It was later noticed that one of the selected problems was satisfiable, and removed, leaving
% 100 problems for the competition.
% Eighty-four of the 100 problems are existing TPTP problems, of which 32 had difficulty rating 
% 1.00, and the average difficulty rating over the 84 was 0.90.
% The remaining 16 problems were all versions of existing TPTP problems that had been modified
% to suit the entrant's system.
% 
% The problems were given in reverse order of desired use, so that the ``easier'' problems were
% used before ``harder'' ones.
% 
%---------------------------------------------------------------------------------------------------
% \subsection{Problems for the LTB Division}
% \label{LTBProblems}
%
% The problems for the LTB division are taken from various sources,
% % (both within and across CASC editions),
% with each problem category being based on one source.
% \THIS{} had only one category,
% % \footnote{%
% % Usually the LTB division has multiple problem categories, but for CASC-J9
% % it was decided to have just one category, which allowed that category
% % to have a very large number of problems.}:
% the VBT category, which used a set of 8900 problems generated by Isabelle's Sledgehammer
% system from the ``van Emde Boas Trees'' session in the AFP.\footnote{%
% \url{https://www.isa-afp.org/entries/Van_Emde_Boas_Trees.html}}
% Recall from Section~\ref{Introduction}, six versions of each problem were provided:
% monomorphic and polymorphic typed first-order versions (TF0 and TF1),
% monomorphic and polymorphic typed extended first-order versions (TX0 and TX1),
% and
% monomorphic and polymorphic typed higher-order versions (TH0 and TH1).
% % Eight versions of each problem were provided -- the six versions used in the competition as
% % described in Section~\ref{Introduction}, and additionally TH0 and TH1 versions that used
% % {\tt \$ite}, {\tt \$let}, or tuples.
% % As some systems were not able to attempt problems with {\tt \$ite}, {\tt \$let}, or tuples, it
% % was decided to exclude those two versions from the competition.
% In order to gauge the problems' difficulties and find proofs that could be used in the training
% data that is provided in the LTB division, five ATP systems (E~2.6.X, lambdaE~22.03.23,
% Leo-III~1.6.6, Vampire~4.6.1, Zipperposition~2.1) were run on all versions of the 8900 problems
% with a 60~s CPU time limit.
% A prerelease of 100 problems with their proofs was evenly sliced out from the 8900 problems and
% distributed to the entrants before the competition so there were no surprises, leaving 8800
% problems.
% % The FOF versions were attempted by CSE\_E 1.0, E 2.2, iProver 2.8, Metis 2.4,
% % Prover9 1109a, SPASS 3.9, and Vampire~4.3.
% % The TF0 versions were attempted by E 2.2, iProverModulo 2.5-0.1, Leo-III 1.3,
% % Vampire~4.3, and Zipperposition 1.4.
% % The TF1 problems were attempted by Leo-III 1.3 and Zipperposition 1.4.
% % The TH0 problems were attempted by agsyHOL 1.0, cocATP 0.2.0, LEO-II 1.7.0,
% % Leo-III 1.3, Satallax 3.3, and Zipperposition 1.4.
% % The TH1 problems were attempted by
% % HOLyHammer 0.21 and Leo-III 1.3.
% % A proof for at least one version was found for 10294 problems, at least one proof for every
% % version was found for 2632 problems, and all the systems solved all the versions of 29 problems.
% The problems were then sorted according to the number of solutions found across all versions, and
% a training set of 800 problems that each had at least two proofs was sliced out as the training
% data to be used in the competition.
% The remaining 8000 problems were used in the competition.
% There were 3517 problems where all versions had at least one proof (from the five ATP systems),
% 6867 problems where at least one version had at least one proof,
% 800 problems where at least one version had no proofs, and
% 490 problems that had no proofs.
% These results provided assurance that the problems were of a reasonable range of difficulty.
% % and the 490 unsolved problems provided a challenge to the systems in the competition
% % (see Section~\ref{LTB}).
% The TF0 problems used had up to 12792 axioms,
% the TF1 problems had up to 12528 axioms,
% the TX0 problems had up to 12579 axioms,
% the TX1 problems had up to 12488 axioms,
% the TH0 problems had up to 10296 axioms,
% and
% the TH1 problems had up to 10078 axioms.
%
% The problems are given to the ATP systems as files in TPTP format, in the natural order of their
% source, i.e., for \THIS{} in the order of their export from the AFP.
% The problem batch was \textit{unordered}, which meant that the ATP systems could attempt as many
% of the versions as they wanted, in any order including in parallel, and could attempt them
% multiple times.
% This provided opportunities for sharing logical and control results between proof attempts.
% A solution to any version counted as a solution to the problem.
%
% % \begin{packed_itemize}
% % \item The AIM2175 problem set contains 1155 test problems, and 1020 training
% %       problems with one solution each.
% %       For the competition, 200 problems were randomly chosen from the 1155.
% %       All the training problems and their solutions were given in the
% %       training set.
% % \item The CML10227 problem set contains 10227 problems.
% %       Proofs were found by Vampire~4.0 for 2017 of them by careful
% %       preselection of axioms.
% %       The 2017 proofs were used to construct ``easy problems'' with only the
% %       axioms used in the proofs, which were attempted in non-batch mode by
% %       pre-\THIS{} versions of E, iProver, and Vampire, with a 120~s CPU time
% %       limit.
% % %      (Those ATP systems were chosen as they were the underlying systems
% % %      entered into the (non-demonstration) division.)
% %       Only 440 of these easy problems were solved by any of the three
% %       systems.
% %       One hundred of the 440 solved problems and 100 of the 1577 unsolved
% %       problems were randomly selected to be used in the competition, but
% %       with their full original axiomatizations.
% % % (i.e., not with just the axioms extracted to form the easy problems).
% %       For the training set, 1000 problems (and their Vampire~4.0 solutions)
% %       were selected from the remaining 1817 unselected problems.
% % \item The HH7150 problem set contains 7150 problems.
% %       This problem set was also used in \LAST{}, and same criteria were
% %       used for problem selection:
% %       140 problems that were solvable in 60~s by at least one of the \LAST{}
% %       non-LTB versions of the systems, and 60 problems that were not solvable
% %       in 60~s.
% %       For the training set, 1000 problems (and their solutions) were selected
% %       from those that were solvable in 60~s but not selected for the
% %       competition.
% % \end{packed_itemize}

% For the LTB division this preprocessing was limited to randomly
% reordering the formulae, so that the formulae in the training set were
% the same as in the competition problems.
% In the non-batch divisions
% The problems are given
% individually as command line parameters.
% In the LTB division the problems for each problem category are listed in a
% \textit{batch specification file}, containing global data lines and one or more
% batch specifications.
% The global data lines specify the problem category, and for CASC-24
% provided the name of the directory containing the training data.
% A batch specification consists of a configuration section, a list of the
% common axiom files \texttt{include}d in every problem, and a list of the
% problems in the batch.

%--------------------------------------------------------------------------------------------------
\subsection{Time Limits}
\label{TimeLimits}

% In the THF, FOF, FNT, and UEQ divisions 
A wall clock time limit is imposed on each system's execution.
The minimal time limit is 120~s.
% , except for the ICU division where the minimal time limit is 300~s.
The maximal time limit is constrained by the same factors that constrain the numbers of problems 
that are used, taking into account the phenomenon that ATP systems solve most problems quickly 
and very few slowly~\cite{SS01,Sut24}.
(This phenomenon is evident in the performance plots from the competition.\footnote{%
\href{https://tptp.org/CASC/J13/WWWFiles/ResultsPlots.html}{\tt tptp.org/CASC/J13/WWWFiles/ResultsPlots.html}})
% and from the ratios of total times taken to solved times given in Section~\ref{Computers} --
% 1618 vs. 185 CPU hours, and 453 vs. 34 wall clock hours.
The time limit is chosen value within the range allowed according to the judgement of the
organizer, and is announced at the competition.
In \THIS{} a 180~s wall clock time limit was imposed on each system's execution.
No CPU time limits are imposed, so it can be advantageous to use all the cores on the CPU.
% when a wall clock time limit is used.

% \footnote{%
% CPU limits emphasize underlying algorithmic/search quality, while wall
% clock limits emphasize the user experience.
% For many years, up to CASC-27 in 2019, CPU limits were used.
% At that time it was decided to change to wall clock limits to promote the
% use of the multiple cores that had become commonly available.
% There will be a time in the future when development of multi-core techniques
% have matured, and CASC might then go back to CPU limits.}
% A wall clock time limit of 360~s was also imposed to limit very high memory
% usage that causes swapping.

% In the SLH division a CPU time limit is imposed for each problem.
% The limit is between 15~s and 90~s, which is the range of CPU time that can be usefully
% allocated for a proof attempt in the Sledgehammer context.\footnote{%
% According to a personal communication from Jasmin Blanchette, and he should know.}
% The time limit is chosen within the range allowed according to the judgement of the organizer,
% and is announced at the competition.
% In \THIS{} a 15~s CPU time limit was imposed for each problem.
% 
% In the ICU division a wall clock time limit is imposed for each problem.
% The limit is between 300~s and 600~s, which is a range that gives the systems sufficient time
% (4800~s CPU time on the octa-core CPUs) to attempt the difficult problems submitted.
% The time limit is chosen within the range allowed according to the judgement of the organizer,
% and is announced at the competition.
% In \THIS{} a 480~s wall clock time limit was imposed for each problem.

% In the LTB division a 172800~s wall clock time limit was imposed for
% the one batch in
% the problem batch, giving an average of 21.6~s per problem.
% There was no wall clock time limit for each problem, and no CPU time limits.
% Time spent before starting the first problem in the batch, e.g., tuning
% on the training data and pre-loading common axioms, and time spent
% between ending a problem and starting the next, e.g., learning from previous
% proofs, is not part of the time taken on problems.
% However, the time taken on such tasks is part of the overall time for the
% batch.

%---------------------------------------------------------------------------------------------------
\section{System Registration, Delivery, and Execution}
\label{Registration}

The ATP systems are registered for the competition 
% by the registration deadline, 
using an system registration form,
Systems can be entered at only the division level, and can be entered into more than one division.
Entering many similar variants of a system is deprecated, and entrants may be required to limit
the number of variants that they enter.
Systems that rely essentially on running other ATP systems without adding value are deprecated, and
such systems might be disallowed or moved to the demonstration division.
For each system an entrant must be nominated to handle all issues (e.g., installation and execution
difficulties) arising before, during, and after the competition.
The nominated entrant must formally register for CASC.
It is not necessary for entrants to physically attend the competition.

The ATP systems are delivered to the competition organizer 
% by the system delivery deadline,
as StarExec \smalltt{.tgz} installation packages.
The installation packages can contain only the components necessary for running the system, i.e.,
not including source code, etc.
The competition organizer installs and tests the packages on StarExec.
Source code is delivered separately, under the trusting assumption that the installation package
corresponds to the source code.
% The entrants must submit a \smalltt{.tgz} file of the source code and any files required for
% building the StarExec installation package to the competition organizer.
% by the system delivery deadline.

After the competition all competition division systems' StarExec and source code packages are made
available on the competition web site.
This allows anyone to use the systems on StarExec, and to examine the source code.
Entrants are encouraged to make a public release of their systems after the competition, so
that users can enjoy the latest capabilities.
An open source license is encouraged, to allow the systems to be freely used, modified, and
shared.
Many of the StarExec packages include statically linked binaries that provide further portability
and longevity of the systems.
In the demonstration division the entrant specifies whether or not the source code is placed on
the web site.

% For systems running on entrant supplied computers in the demonstration division, entrants must
% email a \smalltt{.tgz} file containing the source code and any files required for building the
% executable system to the competition organizer by the system delivery deadline.
% In the demonstration division the entrant specifies whether or not the source code is placed on
% the CASC web site.

Execution of the ATP systems is controlled by StarExec.
StarExec copies the systems and problems to the compute nodes before starting execution, and timing
starts when execution starts.
The ATP systems must be fully automatic -- they are executed as black boxes, on one problem
% (in the non-LTB divisions) or one batch (in the LTB division)
at a time.
Any command line parameters have to be the same for all problems in each division.
The ATP systems must be sound, and are tested for soundness before the competition by submitting
non-theorems to the systems in the THF, FOF, and UEQ divisions, and theorems to the systems in the
FNT division.
Claiming to have found a proof of a non-theorem or a disproof of a theorem indicates unsoundness.
If a system fails the soundness testing it must be repaired by the unsoundness repair
deadline or be withdrawn.
Three systems were found to be unsound before \THIS{}, and were repaired in time for the 
competition.
The execution of demonstration division systems is supervised by their entrants.

After the competition the systems are tested over the whole TPTP.
If a system is found to be unsound, but before the competition report is published, and it cannot
be shown that the unsoundness did not manifest itself in the competition, then the system is
disqualified.
After the competition the systems' solutions are checked.
If any of the solutions are found to be unacceptable, but before the competition report is
published, then the ranking is revised.
All disqualifications and ranking changes are explained in the competition report.

In CASC each ATP system is entered under one name, but it is common for the entrant to deliver
multiple installation packages, each specialized to particular divisions in the competition,
e.g., separate packages for the FOF and UEQ divisions.
% In some cases multiple packages are combined into one by adding a wrapper that detects the 
% SPC of the problem, and then calls the appropriate package's entry point.
Requiring the use of a different package depending on the problem type is a coarse level of 
non-automation.
Users of ATP systems can have different intentions that produce distinct tasks for ATP, the
most common of which is theorem proving/unsatisfiability checking vs. countermodel 
finding/satisfiability checking
In \NEXT{} entrants will be restricted to maximally two installation packages, one for each
of these two types of task.
The ATP systems will be responsible for detecting the problem type within each type of task,
and running appropriate components.
This will improve the level of automation in ATP systems.

%--------------------------------------------------------------------------------------------------
\section{System Evaluation}
\label{Evaluation}

CASC ranks the ATP systems at (only) the division level.
For each system, for each problem, four items of data are recorded:
whether or not the problem was solved,
whether or not a solution (proof or model) was output,
the CPU and wall clock times taken to solve the problem and the times taken to output the
solution (taken from the time stamps prepended to each line of the system's \smalltt{stdout}
by StarExec's \smalltt{runsolver} utility),
and
the maximal number of parents in any inference.
The systems are ranked according to the number of problems solved with an acceptable solution 
output.
Ties are broken by the average time taken over problems solved.
% (The models are ``counter-models'' for problems with an unprovable conjecture,
% and ``models'' for satisfiable axiom sets without a conjecture.)
Trophies are awarded to the competition divisions' winners.
A system that is not entered into a division is assumed to perform worse than the entered systems,
for that type of problem -- wimping out is not an option.
% Additionally in CASC-28, Jasmin Blanchette at the Vrije Universiteit Amsterdam
% contributed a travel prize of a ``romantic hotel stay in Amsterdam'' for the
% SLH division, with expenses covered by the Matryoshka project\footnote{%
% \url{https://matryoshka-project.github.io}}~\cite{BF+17}.
In the demonstration division the systems are not ranked, and no trophies are awarded.

The competition panel decides whether or not the systems' solutions are ``acceptable''. 
The criteria include:
\begin{packed_itemize}
\item Proofs should not use a TPTP language that is more expressive than that of the problem.
\item Proofs must be syntactically correct in TPTP format~\cite{Sut07-CSR}.
      Proof syntax is checked using the BNF-based parser~\cite{VS06} and TPTP4X~\cite{Sut10}.
      This can be done in SystemOnTSTP.\footnote{%
\href{https://tptp.org/cgi-bin/SystemOnTSTP}{\tt tptp.org/cgi-bin/SystemOnTSTP}}
\item Proofs must be structurally correct:
      \begin{packed_itemize}
      \item Annotated formulae must be uniquely named.
      \item Proofs must be acyclic.
      \item Proofs must show only relevant inference steps.
      \item Proofs must have formulae from the problem as leaves, and end at the conjecture (for
            axiomatic proofs) or \smalltt{\$false} formulae (for proofs by contradiction,
            e.g.,~CNF refutations and closed tableau).
      \item Inference steps should list only parents that are used in the inference.
      \item Proofs that make \smalltt{assumption}s must propagate and eventually discharge the
            assumptions.
      \end{packed_itemize}
      Proof structure is checked using GDV~\cite{SB05}.
      This can be done in SystemOnTSTP.
\item Translations from one form to another, e.g., FOF to CNF, must be adequately documented.
\item Inference steps must be reasonably fine-grained.
      The maximal number of parents in any inference is considered in this regard.
\item Proofs should graft in subproofs from external systems.
\item An unsatisfiable set of ground instances of clauses is acceptable for establishing the
      unsatisfiability of a set of clauses.
\item Models must be complete, documenting the domain and symbol maps.
      The domain and symbol maps may be specified by explicit ground lists (of mappings), or by
      any clear, terminating algorithm.
      Saturations are acceptable if their models are a subset of the models of the problem formulae.
      Use of the TPTP format for interpretations~\cite{SS+26} is encouraged.
\end{packed_itemize}

In addition to the ranking data, other measures are made and presented in the results:
\begin{packed_itemize}
\item The \textit{state-of-the-art contribution} (SotAC) quantifies the unique abilities of each
      system (excluding the previous year's winners that are earlier versions of competing systems).
      For each problem solved by a system, its SotAC for the problem is the fraction of systems 
      that do not solve the problem.
      A system's overall SotAC is the average over the problems solved but that are not solved 
      by all the systems.
\item The \textit{efficiency} balances the number of problems solved with the time taken.
      It is the average solution rate (the reciprocal of the time taken to solve it) over the 
      problems solved, multiplied by the fraction of problems solved.
      Efficiency is computed for both CPU time and wall clock time, to measure how efficiently the
      systems use one core and multiple cores respectively.
\item The \textit{core usage} measures the extent to which the systems take advantage of multiple
      cores.
      The raw core usage is the ratio of CPU time to wall clock time used, and the average core
      usage is over the problems solved with a raw core usage greater than 1.0.
      The competition ran on octa-core computers, thus the maximal core usage was 8.0.
      The ability to solve problems quickly before multi-core search is used is also of interest.
      The number of problems solved with a raw core usage less than 1.0 is noted as an imprecise
      cutoff for separating proofs that are found in preprocessing on a single core, from proofs
      found later after multi-core search has started. 
      It is understood that some systems use multiple cores from before 1~s.
\item The \textit{granularity} measures the extent to which inference steps combine multiple 
      parent formulae into a single inferred formula.
      The granularity of a proof is the maximal number of parents in an inference, and the 
      overall granularity is the average over the problems solved with acceptable proof output.
      Granularity is a new measure in \THIS{}. 
      It was motivated by the very large numbers of parents listed in individual inferences of 
      \LAST{}, as discussed in the competition report~\cite{Sut26-CASC}.
%\item In the ISA problem category of the LTB division, the {\em number
%      of axiom lists} (output of a list of axioms sufficient for a proof)
%      is counted.
%      A well-formed proof can be used to provide such a list.
%\item In the SMO problem category of the LTB division, the number of
%      {\em questions answered} (output of the bindings for the outermost
%      existentially quantified variables) is counted
%      (see Section~\ref{Output}).
\end{packed_itemize}

% In the ISA problem category of the LTB division, the \textit{axiom accuracy}
% compares the numbers of axioms used in proofs.
% For each problem solved by a system, its problem axiom accuracy is
% the smallest number of axioms used in any system's proof, divided by the
% number of axioms in this system's proof (or 0 if this system does
% not output a proof).
% A system's overall axiom accuracy is its average problem axiom
% accuracy over the problems it solved, multiplied by the fraction of
% problems solved.
% This was the basis for the ISA category prize.
% In the SMO problem category of the LTB division, the \textit{number of questions
% answered}, i.e., output of the bindings for the outermost existentially
% quantified variables, is counted.
% This was the basis for the SMO category prize.

%--------------------------------------------------------------------------------------------------
\section{Results}
\label{Results}

% This section summarizes the results of \THIS{}, and provides commentary.
% For the various divisions and categories, according to their ranking
% criteria (see Section~\ref{Evaluation}),
The result tables below give 
the number of problems solved with an acceptable solution output,
the number of problems solved without considering solution output,
the average wall clock time taken over the problems solved,
the state-of-the-art contribution,
the (micro-)\allowbreak efficiency,
the core usage/the number of problems solved with core usage greater than 1.0,
the proof granularity,
the number of problems uniquely solved,
the number of new problems solved,
and
the number of problems solved in each problem category.
In each column the best value in the competition division is {\bf emboldened}.
The \LAST{} winner is \textit{emphasized} in the demonstration divisions.
Detailed results, including the systems' output files, are available on the \THIS{} web 
site.\footnote{%
\href{https://tptp.org/CASC/30/WWWFiles/Results.html}{\tt tptp.org/CASC/30/WWWFiles/Results.html}} 

The structural checking of proofs in the competition theorem proving divisions 
was new in \LAST{}, and many systems' proofs had structural flaws.
Those results were expected to provide motivation for developers of systems in \THIS{} to ensure 
that their systems' proofs are acceptable.
Some developers knew in advance that their systems' proofs would not be acceptable, and their
numbers of problems solved with an acceptable solution output is set to ``-'' rather than 0
in the results tables.
Discussion of the progress in proof output from \LAST{} to \THIS{}, in terms of proof structure
and inference granularity, is given in Section~\ref{Acceptable}.

%--------------------------------------------------------------------------------------------------
\subsection{The THF Division}
\label{THF}

Table~\ref{THFSummary} summarizes the results of the THF division.
Vampire~5.0.1 won the division, but solved less problems than the \LAST{} winner Vampire~5.0,
which in \LAST{} had solved less problems than the preceding CASC-J12 winner!
The versions of Vampire up to and including 5.0.1 have separate binaries for first-order and 
higher-order logics.
The higher-order components of Vampire are being incrementally merged into the main branch
of development, which is clearly producing some growing pains.
Zipperposition ran the same code as it has since CASC-29 in 2023, and its proofs that were 
acceptable in the past are now unacceptable due to the stricter requirements.
The flaws in Zipperposition's proofs are quite simple, e.g., many of them have duplicate formulae 
names, in some cases duplicate formulae, and others give the SZS \smalltt{status(esa)} when 
negating the conjecture (it should be \smalltt{status(cth)}).
These are flaws that should be easily fixed, so that Zipperposition is upgraded for \NEXT{}.
The 13 missing proofs from Leo-III are interesting: they come from Leo-III using bracketed
quantifiers as terms that are applied, e.g., \smalltt{((!)\,@\,\ldots)}, instead of the
combinator equivalents expected in the TPTP language, e.g., \smalltt{(!!\,@\,\ldots)}.

\begin{table*}[tbh]
\caption{THF division results}
\label{THFSummary}
\begin{tabular}{lrrrrrr@{}l@{}rrrrrr}
\hline
ATP System             & Solns & Solved & Avg  & Sot  & WC       & \multicolumn{2}{r}{Core} & TNE & TEQ &Uniq & New & Gr'ty \\
                       & /400  & /400   & WC   & AC   & $\mu$Eff & \multicolumn{2}{r}{usage}&/100 &/300 & /31 & /31 &       \\
\hline
Vampire~5.0.1          &  \bb{340}  &  \bb{352}   &  \bb{6.5} & \bb{0.32} & \bb{615}      & 3.9       & {\tiny /212} &  67 & \bb{285} &   3 &  \bb{31} & \bb{10.1}  \\
E~3.5.1                &  323  &  324   & 10.7 & 0.29 & 593      & 4.2       & {\tiny /129} &  \bb{73} & 251 &  \bb{23} &  28 & 26.4  \\
Leo-III~1.8.0          &  137  &  160   & 14.4 & 0.07 &  83      & 2.5       & {\tiny /160} &  38 & 122 &   1 &   7 & 53.3  \\
LEO-II~1.7.0           &   67  &   67   & 11.0 & 0.01 & 110      & 0.0       & {\tiny /0  } &  26 &  41 &   0 &   0 & 84.8  \\
Zipperp'n~2.1.9999     &    0  &  347   & 16.7 & 0.31 & 334      & \bb{5.7}       & {\tiny \bb{/268}} &  67 & 280 &   1 &  29 &    -  \\
\hline
\multicolumn{13}{c}{Demonstration division} \\
\hline
\textit{Vampire~5.0}   &  344  &  354   &  6.8 &    - & 576      & 4.2       & {\tiny /238} &  70 & 284 &   3 &  31 & 14.8  \\
\hline
\end{tabular}
\end{table*}

Zipperposition had the highest average time taken, but made the best use of the multiple cores.
E solved the most problems before starting to use multiple cores, and had good core usage
Leo-III used multiple cores for all the problems it solved.
E solved the most TNE problems, largely because it uniquely solved 19 graph theory 
(\smalltt{GRA}) problems.
These problems have been in the TPTP problem library since release v3.6.0 from 2008.
E's developer suggested that the progress was simply due to several bugs being fixed in the 
higher-order code.
The new problems were mostly solved by the non-Leo systems.

The granularity figures are not extreme, with the Leos having the highest values.
Analysis revealed that their higher granularities came from inferences that were to the E system.
The developer said: ``LEO-II does exactly this: Just list all of the clauses that were sent to 
E as inference parents. \ldots Leo-III does some naive filtering {\em [of the clauses sent to E]}, 
but it does not filter the parents wrt. actual usage by E. This is ongoing work, but will 
definitely lead to more realistic/lower number of inference parents on average.''
The largest contributor to E's granularity was deeply nested \smalltt{apply\_def} steps, in 
which definitions were sequentially applied in a single inference.

The individual problem results (without considering proof output) show that
11 problems were unsolved,
53 problems were solved by all the systems, and
31 problems were solved by only one system.
% (6 problems were solved by only the two versions
% of Vampire, and are counted as unique solutions for Vampire 5.0.1).
Of the 31 unique solutions,
23 were by E (including the 19 \smalltt{GRA} problems),
3 were by Vampire~5.0.1,
3 were by Vampire~5.0,
1 by Leo-III,
and
1 by Zipperposition.
A portfolio (sharing the time between several systems) of Vampire~5.0 and E, with 90~s
allocated to each, would solve 383 problems.

%--------------------------------------------------------------------------------------------------
% \subsection{The THN Division}
% \label{THN}

% Table~\ref{THNSummary} summarizes the results of the THN division.

% \begin{table}[tbh]
% \caption{THN division results}
% \label{THNSummary}
% \begin{tabular}{lrrrrrrrr}
% \hline
% ATP System            & THN & Avg  &Mdls & Sot  & \multicolumn{1}{c}{$\mu$} & New & TNN & TNQ \\
%                       & /200& CPU  & out & AC   & Eff. & /5  &/25  &/175 \\
% \hline
% Nitpick 2015          & 200 &  7.9 &\tic & 0.70 & 130  &   5 &  25 & 175  \\
% Refute 2015           &  74 & 24.3 & \xx & 0.49 &  29  &   2 &  25 &  49  \\
% Satallax 2.8          &  49 &  0.1 & \xx & 0.48 & 242  &   2 &   5 &  44  \\
% \hline
% \multicolumn{9}{c}{Demonstration division} \\
% \hline
% Isabelle-HOT 2012 & 166 & 88.4 &     & 0.28 &  28  &  12 &  34 & 132 \\
% \hline
% \end{tabular}
% \end{table}

%--------------------------------------------------------------------------------------------------
% \subsection{The TFA Division}
% \label{TFA}
% 
% Table~\ref{TFASummary} summarizes the results of the TFA division.
% The winner was Vampire~5.0
% Vampire's arithmetic reasoning capabilities have recently been improved by the ALASCA approach 
% to reasoning over quantified linear arithmetic~\cite{KK+23}, with further extensions by the VIRAS
% technique for quantifier elimination in integer-real arithmetic~\cite{SKK24}.
% 
% \begin{table*}[tbh]
% \caption{TFA division results}
% \label{TFASummary}
% \begin{tabular}{lrrrrrr@{}l@{}rrrrr}
% \hline
% ATP System             & Solns & Solved & Avg  & Sot  & WC       & \multicolumn{2}{r}{Core} & TFI & TFE & New \\
%                        & /150  & /150   & WC   & AC   & $\mu$Eff & \multicolumn{2}{r}{usage}& /100& /50 & /7  \\
% \hline
% Vampire~5.0            & 127  & 131   &  4.4& 0.40& 695     & 4.7      & {\tiny /72}  & 97& 37& 6  \\
% iProver 3.9.3          &  68   &  74    & 12.0 & 0.10 & 162      & 4.1       & {\tiny /68}  &  49 &  25 & 4   \\
% \hline
% \multicolumn{9}{c}{Demonstration division} \\
% \hline
% \textit{Vampire 4.9}   &   -   & 123    &  4.6 &    - & 592      & 4.8       & {\tiny /80}  &  87 &  36 & 7   \\
% cvc5 1.3.0             &   -   & 108    & 16.4 & 0.18 & 349      & 0.0       & {\tiny /0}   &  75 &  33 & 5   \\
% \hline
% \end{tabular}
% \end{table*}
% 
% The SotAC, efficiency, category, and new problem rankings are aligned with the overall ranking.
% There were 7 new problems in the division, all in the TFI problem category, i.e., using integers.
% Four of the new problems are software creation problems (the TPTP {\tt SWC} domain) from the
% B~formal method~\cite{Deh22-Zenodo}, which were all solved by all the systems.
% Two of the problems are common sense reasoning problems (the TPTP {\tt CSR} domain) that encode
% calculus situations~\cite{MS05}, both of which were solved by Vampire~4.9 but only one by 
% Vampire~5.0.
% The seventh problem is an arithmetic problem (the TPTP {\tt ARI} domain) that was solved by
% only the two Vampires.
% 
% The individual problem results (without considering proof output) show that
% 15 problems were unsolved,
% 65 problems were solved by all the systems, and
% 23 problems were solved by only one system (16 problems were solved by only the two versions
% of Vampire, and are counted as unique solutions for Vampire~5.0).
% Of the 23 unique solutions,
% 19 were by Vampire~5.0,
% 3 by cvc5,
% and
% 1 by Vampire 4.9
% The dominance of Vampire~5.0 means that only a little can be gained from a portfolio of the 
% systems -- a combination of the two Vampires with cvc5, with 40~s allocated to each, would solve 
% 133 problems.
% 
% The \LAST{} report said:
% ``Over the years the TFA division has struggled to attract entrants, with the same few systems
% being entered each year~\ldots~CASC will continue to have a TFA division in order to encourage 
% development of ATP in this area.''
% Sadly there has been no change to the situation, and the TFA division will not be run in
% \NEXT{} (which will be somewhat different anyway, as explained in Section~\ref{Conclusion}).

%--------------------------------------------------------------------------------------------------
% \subsection{The TFN Division}
% \label{TFN}
% 
% Table~\ref{TFNSummary} summarizes the results of the TFN division.
% Vampire~5.0 solved the most problems, iProver~3.9.3 came in second, and cvc5 third.
% All three systems output acceptable models for all the problems they solved.
% This is in contrast to \LAST{} where Vampire~4.9 solved the most problems but none of its
% models were considered acceptable, iProver~3.9 came second but output acceptable models for
% only 67 of the 85 problems it solved, and cvc5 came third but output no acceptable models.
% This matter was closely examined after \LAST{}, and the developers of all three systems
% admitted to the deficiencies in their systems' model output.
% It is pleasing to see the positive impact that had, leading to improved model output from all 
% three systems in \THIS{}.
% 
% \begin{table}[tbh]
% \caption{TFN division results}
% \label{TFNSummary}
% \begin{tabular}{lrrrrrr@{}l@{}r}
% \hline
% ATP System           & Solns & Solved & Avg  & Sot  & WC       & \multicolumn{2}{r}{Core} & New \\
%                      & /150  & /150   & WC   & AC   & $\mu$Eff & \multicolumn{2}{r}{usage}& /2      \\
% \hline
% Vampire 5.0          & 110  & 110   & 1.3 & 0.27& 588     & 4.8       & {\tiny /71}  & 1     \\
% iProver 3.9.3        &  86   &  86    & 23.0 & 0.12 & 159      & 5.8      & {\tiny /60}  & 1     \\
% cvc5 1.3.0           &  56   &  56    & 28.5 & 0.05 & 193      & 0.0       & {\tiny  /0}  & 1     \\
% \hline
% \multicolumn{8}{c}{Demonstration division} \\
% \hline
% \textit{iProver 3.9} &  84   &  84    & 21.5 & 0.11 & 193      & 6.3       & {\tiny /52}  & 1      \\
% \hline
% \end{tabular}
% \end{table}
% 
% As noted in Section~\ref{TPTPProblems}, 68 problems with rating 0.00 and 49 problems
% with rating 1.00 were made eligible for the TFN division.
% Of those, 59 problems of rating 0.00 and 44 problems of rating 1.00 were selected.
% Of the 59 problems with rating 0.00, 41 were solved by all the systems, and all the problems were 
% solved by Vampire~5.0.
% Of the 44 problems with rating 1.00, 39 were unsolved, and the other 5 were solved by only 
% Vampire~5.0.
% One might expect all the problems with rating 0.00 and none of the problems with rating 1.00
% to be solved, but this is not the case because the ratings were calculated using earlier versions
% of the systems. 
% Progress in ATP is not monotonic.
% 
% The individual problem results (without considering model output) show that
% 40 problems were unsolved,
% 46 problems were solved by all the systems, and
% 14 problems were solved by only one system (all 14 problems were solved by only Vampire~5.0).
% % Of the 14 unique solutions,
% % 11 were by Vampire,
% % and
% % 1 was by iProver.
% A portfolio approach cannot improve on the individual systems' results.
% Of the 14 problems solved by only Vampire~5.0, 
% 8 are hardware verification (the TPTP {\tt HWV} domain) problems, 
% 3 are interactive theorem proving (the TPTP {\tt ITP} domain) problems, 
% 2 are software verification (the TPTP {\tt SWW} domain) problems, and
% 1 is a logic calculus (the TPTP {\tt LCL} domain) problem, 
% One had rating 0.00, 7 had rating 0.33, 4 had rating 0.67, and 2 had rating 1.00.
% The common feature of these problems, with the exception of the logic calculus problem, is
% that they have quite large numbers of formulae, ranging from 254 to 7803 formulae.
% Twelve of them were solved by finding a finite model, and two by saturation.
% Vampire's strategy scheduling allows it to try multiple approaches to establishing satisfiability,
% and the ability to find finite models gave it a clear advantage.

%--------------------------------------------------------------------------------------------------
\subsection{The FOF Division}
\label{FOF}

Table~\ref{FOFSummary} summarizes the results of the FOF division.
As in \LAST{}, the Vampires dominated the division, but this year there is clear progress
from the previous version to the new version of the system.
The developers explained: ``most of the improvement between 5.0 and 5.0.1 comes from the addition 
of the neurally guided clause-selection strategies~\cite{Sud25} \ldots~Since we solve many hard 
(previously unsolved) problems with the neurally guided strategies, I think a key was to submit 
the new version of Vampire early, for the problem rating, so that these problems would become 
eligible for CASC and could showcase the improvement.''.
Table~\ref{FOFSummary} and the results plot\footnote{%
\url{https://tptp.org/CASC/J13/WWWFiles/ResultsPlots.html\#FOFProblems}} show clear grouping
of the systems, but the grouping has changed since the last two CASC editions.
In \LAST{} there were three groups: the Vampires, a middle group of six systems including E, 
iProver, Drodi, and Zipperposition, then a lower group of four systems.
It was noted in the \LAST{} report that the grouping was less pronounced than in CASC-J12.
This year the grouping is even less pronounced, but three groups are evident:
the Vampires are out front as usual, then there is a group of CSI++, E, and iProver,
followed by a middle group of Drodi, LisaST, Prover9~2026-6A, and Zipperposition, then a lower 
group of FindProof, mrs, Prover9---1109a, and SUPr, and finally the bottom group.
% of Connect++, SATResetCoP, SPASS-SCL (noting that SPASS-SCL cannot attempt problems with 
% equality), and VIP.
It is not easily possible to find common features of the systems in each group, other than
years of developmental experience.
The top groups include systems that have many years of active development and participation in 
CASC, while the lower groups include newcomers.
CASC clearly provides opportunities for system developers to learn from others!

In \LAST{} a key differentiator between the FOF division systems was their abilities to generate 
acceptable proofs, with a notable difference between the ranking by number of acceptable proofs 
output and the ranking by the number of problems solved.
It is pleasing to see a better correlation between the two rankings in \THIS{}, indicating that
lessons were learned from the \LAST{} -- more analysis of the improvement is in 
Section~\ref{Acceptable}.
However, there is still room for further improvement.
This year the unacceptable proofs often had simple syntax mistakes, that can be easily fixed, 
e.g.,~illegal formula names, missing brackets around binary formulae.
Some other flaws were deeper, e.g., missing formulae.
Vampire~5.0.1 had one proof with a syntax error.
LisaST suffered the most, with a mixture of syntax mistakes and deeper flaws.
CSI - 5 found by own inferences (grep tabproof:). Were any not solved by E? How many others helped E?
\smalltt{SEU269+1} solved by CSI++, but no SZS status line. Also solved by E.
\smalltt{SYN986+1.004} not solved by E alone, but solved by E in CSI++..
ASKED THEM.
In \NEXT{} the requirements for proofs to be considered ``acceptable'' will be tightened, as 
explained in Section~\ref{Acceptable}.

A highlight of \THIS{} was the number of newcomers to the competition.
In the FOF division there were four absolute newcomers: FindProof, SUPr, mrs, and VIP.
Additionally there were four new variants of existing systems: CSI++, Prover9~2026-6A, LisaST, and
SATResetCoP.
It was particularly pleasing to see the new release of Prover9, which earned it the ``Honourable 
Mention'' medal for CASC in the FLoC Olympic games.
The old Prover9~1109a has been the fixed point in the CASC FOF division since CASC-J6 in 2013 
(see Section~\ref{DivisionsSystems}). 
The old version was clearly outperformed by the new version.
The developers explained that the main reason for the new version's strong performance is
a sliding window scheduler that rotates through 96 strategies that had been found using machine 
learning.
It is also noteworthy that the new Prover9 outputs proofs in TPTP format.

\begin{table*}[tbh]
\caption{FOF division results}
\label{FOFSummary}
\begin{tabular}{lrrrrrr@{}l@{}rrrrrr}
\hline
ATP System             & Solns & Solved & Avg  & Sot & WC       & \multicolumn{2}{r}{Core} & FNE & FEQ  &Uniq & New & Gr'ty \\
                       & /400  & /400   & WC   & AC  & $\mu$Eff & \multicolumn{2}{r}{usage}&/100 &/300  & /48 & /79 &       \\
\hline
Vampire~5.0.1          &  \bb{380}  &  \bb{381}   &  \bb{9.5} & \bb{0.58} & 309     & 3.5       & {\tiny /346} &   \bb{94} & \bb{287} &  \bb{45} &  \bb{71} &  6.0  \\
E~3.5.1                &  264  &  264   & 10.6 & 0.32 & \bb{419}     & 4.9       & {\tiny /142} &   75 & 189 &   0 &  55 &  4.3  \\
CSI++~1.0              &  264  &  264   & 14.1 & 0.32 & 406     & 4.6       & {\tiny /139} &   76 & 188 &   0 &  55 &  4.4  \\
iProver~3.9.4          &  222  &  250   & 23.0 & 0.30 & 176     & 4.0       & {\tiny /231} &   64 & 186 &   0 &  54 &265.6  \\
Drodi~4.1.1            &  214  &  215   & 20.3 & 0.24 & 312     & 6.0       & {\tiny /126} &   60 & 155 &   0 &  40 & 10.3  \\
Prover9~2026-6A        &  139  &  143   & 13.0 & 0.13 & 208     & 5.3       & {\tiny /98 } &   50 &  93 &   0 &  25 &  \bb{2.0}  \\
LisaST~0.9             &  127  &  199   & 18.5 & 0.21 &  74     & 5.1       & {\tiny /192} &   54 & 145 &   0 &  52 &  \bb{2.0}  \\
FindProof~0.1          &   91  &   91   & 19.8 & 0.07 &  46     & \bb{6.4}       & {\tiny \bb{/91} } &   26 &  65 &   0 &  18 & 13.3  \\
SUPr~1.0               &   73  &   73   & 16.7 & 0.05 &  92     & 4.0       & {\tiny /54 } &   22 &  51 &   0 &  21 &  3.7  \\
mrs~0.2.0              &   71  &   79   & 12.6 & 0.05 & 106     & 5.1       & {\tiny /59 } &   23 &  56 &   0 &  18 &  2.5  \\
VIP~1.718              &   32  &   32   & 26.1 & 0.01 &  34     & 0.0       & {\tiny /0  } &   13 &  19 &   0 &   8 &  \bb{2.0}  \\
Connect++~0.7.2        &   24  &   64   & 60.0 & 0.04 &  22     & 0.0       & {\tiny /0  } &   27 &  37 &   0 &  13 &  \bb{2.0}  \\
SATResetCoP~1.0        &   20  &   20   & 70.5 & 0.01 &   3     & 0.0       & {\tiny /0  } &   3A &  17 &   0 &   9 &456.0  \\
SPASS-SCL~0.1.1        &   14  &   14   & 21.9 & 0.01 &  12     & 0.0       & {\tiny /0  } &   14 &   0 &   0 &   0 &  4.7  \\
Zipperp'n 2.1.9999     &    0  &  192   & 19.6 & 0.18 & 196     & 4.9       & {\tiny /143} &   46 & 146 &   0 &  47 &    -  \\
\hline
\multicolumn{13}{c}{Demonstration division} \\
\hline
\textit{Vampire~5.0}   &  328  &  328   &  8.1 &    - & 538     & 3.9       & {\tiny /236} &   80 & 248 &   3 &  70 & 13.9  \\
Prover9~1109a          &    -  &   92   & 14.7 &    - & 134     & 0.0       & {\tiny /0  } &   23 &  69 &   0 &  22 &  2.0  \\
\hline
\end{tabular}
\end{table*}

Vampire~5.0.1's high SotAC correlates to its high number of unique solutions.
E's higher efficiency comes from solving most problems very quickly, with relatively fewer being 
solved slowly compared to the Vampires.
This can be seen from E having a sharper upwards swing than the Vampires in the results
plot.
mrs also had a high efficiency relative to its ranking.
FindProof had the highest core usage, using multiple cores in all the problems it solved.
LisaST used multiple cores for most of the problems it solved, which is the intention of
the system design as the developer explained: ``there are 8 built-in strategies, and it 
always runs the 8 in parallel''.
Drodi also had high core usage, and solved 89 problems with a raw cores usage less than 1.
The developer of Drodi provided analysis of the first five and last five problems solved in 
less than 1~s, revealing that eight of those 10 proofs were found during multi-core search.
He explained: ``These results suggest that more than one core is used in most of the problems 
solved in under 1 second and in some cases even a significant number of cores are used.''.
ASKED OSCAR FOR MORE.
Several of the systems, in particular Connect++, SATResetCoP, and VIP, had high average wall 
clock times and no core usage.
The developers of the systems said respectively: ``There could be a few possible reasons for 
this but without some profiling I'd have a hard time making a concrete suggestion.'', ``I bet that 
the issue with SATResetCoP is the Python backend; importing all the libraries needed for the 
initial run is slow.'', and ``WAITING TO HEAR FROM ILIES''.
The expected benefit of using multiple cores to reduce wall clock time taken, as noted in
Section~\ref{TimeLimits}, is evident here.

The granularity values are reasonable for most of the systems.
The best came from Connect++, LisaST, Prover9~2026-6A, and VIP, which all had a granularity of 2.0. 
That means those systems' inferences had an average maximum of two parents, making the proofs as 
fine-grained as can be expected.
At the other end of the scale, iProver and SATResetCoP had high granularity values.
iProver had eight proofs with a granularity of over 1000, three of which were over 10000, with
the maximum being 12651.
As was the case in \LAST{}, iProver's use of Z3 finish refutations is the common culprit,
seen in iProver's inference rule \smalltt{smt\_impl\_just}.
SATResetCoP's proof structure always contains a formulae that is the conjunction of the axioms
and negated conjecture, i.e., the granularity of the proof is the size of the problem.
The problem \smalltt{CSR058+3} has 8006 formulae, and thus the granularity of SATResetCoP's proof 
is also 8006.
Outliers such as these contributed to SATResetCoP's higher granularity value.
The developer said (about these formulae): ``This is generated by the certificate generator, and 
not by SATResetCoP instead. It might be not crucial for the final output''.
% , average 531, median 44.
% ASKED ILIES ABOUT THIS BROKEN PROOF
% (only VIP has one proof with a maximal of one parent
% per inference, in the problem \smalltt{GEO498+1} whose conjecture is already effectively in CNF,
% and was solved by VIP with a single rewrite step using that one formula)

The performances in the two problem categories are well aligned with the overall performance,
except for LisaST's stronger performance on problems with equality, and SPASS-SCL's zero 
performance on problems with equality.
The new problems provided consistent challenges for the systems, with about 20\% of the problems
being new problems, and for most of the systems about 20\% of the problems solved were new 
problems.
The exception is Prover9~2026-6A, for which 17.5\% of the problems solved were new problems.
Careful analysis by the developers concluded that this is: ``a non-significant point estimate 
on a 79-problem sample, not evidence of tuning''.
Only 2 of the 45 problems uniquely solved by Vampire~5.0.1 were new problems.

The individual problem results (without considering proof output) show that
14 problems were unsolved,
no problems were solved by all the systems, and
48 problems were solved by only one system.
% (19 problems were solved by only the two versions of
% Vampire, and are counted as unique solutions for Vampire~5.0).
Of the 48 unique solutions,
45 were by Vampire~5.0.1,
and
3 by Vampire~5.0.
That means that between them, the two Vampires solved a superset of the union of the problems 
solved by other the systems.

% A portfolio of the two Vampires, cvc5, and CSI\_Enigma, with 60~s allocated to each, would 
% solve 484 problems.

%--------------------------------------------------------------------------------------------------
\subsection{The FNT Division}
\label{FNT}

The FNT division return to \THIS{} after two years on hiatus.
The winner from CASC-29 in 2003 was entered as the previous winner.
Table~\ref{FNTSummary} summarizes the results.
Vampire~5.0.1 came out ahead of the other systems, by a large margin, surpassing Vampire~4.8
that had come out way ahead in CASC-29.
The key to the Vampires' success was strategy scheduling that tried to saturate the clause
set, and alternatively find a finite model.
The only other system that produced two types of models was iProver, which output saturations 
and models expressed in a term algebra.
The ratios of saturations to (finite) models output was as follows:
Vampire~5.0.1 - 59:84, 
Vampire~4.8 - 36:91,
iProver - 6:65,
FMB4J 0:51,
Mace4 0:44,
E - 42:0,
and
Drodi - 41:0.

In contrast to the three theorem proving divisions, there was no checking of the models output
in the FNT division.
Post-competition analysis showed that all the competition systems' models passed TPTP syntax 
checking, except for iProver in which some models had duplicate formula names and unquantified 
variables -- these simple errors are easily rectified.
Some of Vampire~4.8's models also had simple errors, but that's history.
% iProver 64 with syntax errors, 42 duplicate formula name, 22 unquantified variables. 
% Vampire 4.8 92 with syntax errors, simple redeclaration of \$i. 
Syntax checking of models will be added in \NEXT{}, as explained in Section~\ref{Acceptable}.
For one problem Mace4 found a model but did not complete the model output by the time limit.

\begin{table*}[tbh]
\caption{FNT division results}
\label{FNTSummary}
\begin{tabular}{lrrrrrr@{}l@{}rrrr}
\hline
ATP System             & Solns & Solved & Avg  &  Sot & WC       & \multicolumn{2}{r}{Core} & FNN & FNQ &Uniq & New  \\
                       & /150  & /150   & WC   &  AC  & $\mu$Eff & \multicolumn{2}{r}{usage}& /50 &/100 & /16 & /10  \\
\hline
Vampire~5.0.1          &  \bb{143}  &  \bb{143}   & 11.9 & \bb{0.58} & 434      & 4.7       & {\tiny /87 } &  \bb{45} &  \bb{98} &  \bb{14} &  \bb{10}  \\
iProver~3.9.4          &   71  &   71   & 25.5 & 0.19 & 176      & 3.3       & {\tiny /53 } &  15 &  56 &   0 &   8  \\
FMB4J~0.1              &   51  &   51   &  9.3 & 0.11 & 241      & 1.3       & {\tiny /33 } &   8 &  43 &   0 &   7  \\
Mace4~2026-6A          &   44  &   45   &  9.0 & 0.10 & 229      & 4.2       & {\tiny /12 } &   8 &  37 &   0 &   8  \\
E~3.5.1                &   42  &   42   &  9.4 & 0.12 & 151      & 4.3       & {\tiny /22 } &  17 &  25 &   0 &   2  \\
Drodi~4.1.1            &   41  &   41   &  6.1 & 0.09 & 175      & \bb{4.9}       & {\tiny \bb{/17} } &   7 &  34 &   0 &   2  \\
Connect++~0.7.2        &    -  &   11   &  \bb{0.4} & 0.01 &  73      & 0.0       & {\tiny /0  } &   3 &   8 &   0 &   0  \\
\hline
\multicolumn{12}{c}{Demonstration division} \\
\hline
\textit{Vampire~4.8}   &  127  &  131   & 17.4 &    - & 411      & 5.6       & {\tiny /82 } &  45 &  86 &  2  &  10  \\
\hline
\end{tabular}
\end{table*}

Vampire~5.0.1 dominated the division across all performance measures.
Drodi stood out as having a low average CPU time coupled with a high core usage -- the insights
from Section~\ref{FOF} apply again here.
The problem category rankings reflect the overall ranking except for E's better performance
on FNN problems.
This is different from CASC-29 in which E~3.1 performed relatively better on FNQ problems.
The developer of E ascribed the change to: ``just random luck'' of the problem selection.

The individual problem results show that
4 problems were unsolved,
2 problems were solved by all the systems, and
16 problems were solved by only one system.
% (21 problems were solved by only the two versions of
% Vampire, and are counted as a unique solutions for Vampire~4.8).
Of the 16 unique solutions,
14 were by Vampire~5.0.1, 
and 2 by Vampire~4.8.
As is the case in the FOF division, between them, the two Vampires solved a superset of union of 
the problems solved by the other systems.

% A portfolio approach does not help much -- a combination of Vampire~4.8, iProver, and cvc5,
% with 60~s allocated to each, would solve just 146 problems.

%--------------------------------------------------------------------------------------------------
% \subsection{The EPR Division}
% \label{EPR}
% 
% The EPR division returned from hiatus in \THIS{}.
% The most recent prior winner was Vampire~4.4 from CASC-27 in 2019. 
% Unfortunately Vampire~4.4's StarExec package has been lost, and was thus not available for the 
% demonstration division.
% A comparison with iProver is a reasonably meaningful comparison system, as iProver dominated the 
% division from CASC-J4 in 2008 through CASC-J9 in 2018.
% Recall also that the EPR division does not require solution output (see Section~\ref{Evaluation}).
% With that context, Table~\ref{EPRSummary} summarizes the results of the EPR division.
% Vampire~5.0 solved the most problems, outperforming iProver.
% 
% \begin{table*}[tbh]
% \caption{EPR division results}
% \label{EPRSummary}
% \begin{tabular}{lrrrrr@{}l@{}rr}
% \hline
% ATP System             &Solved & Avg  & Sot  &  WC      & \multicolumn{2}{r}{Core} & EPU  & EPS  \\
%                        &  /200 & WC   & AC   & $\mu$Eff & \multicolumn{2}{r}{usage}& /100 & /100 \\
% \hline
% Vampire~5.0            &   186&  9.3 & 0.50& 494     & 5.9      & {\tiny /117} &  96 &  90 \\
% iProver 3.9.3          &   168 & 18.7 & 0.41 & 202      & 4.5       & {\tiny /142} &  85  &  83  \\
% CSI\_Enigma 1.0.6      &    98 & 10.2 & 0.09 & 110      & 3.7       & {\tiny  /98} &  32  &  66  \\
% E 3.3.0                &    88 & 19.1 & 0.09 & 294      & 3.5       & {\tiny  /35} &  29  &  59  \\
% Drodi 3.6.0            &    80 &  8.5 & 0.04 & 300      & 6.8      & {\tiny  /27} &  25  &  55  \\
% SPASS-SCL 0.1          &    64 &  4.5& 0.02 & 272      & 3.2       & {\tiny  /21} &  11  &  53  \\
% \hline
% \multicolumn{8}{c}{Demonstration division} \\
% \hline
% No systems             &       &      &      &          &           &              &      &     \\
% \hline
% \end{tabular}
% \end{table*}
% 
% Vampire had a very high SotAC, thanks to it's ability to solve problems that no other systems
% solved (see below).
% Vampire also had the highest efficiency, followed by Drodi.
% Drodi had the highest core usage.
% E suffered from the bug explained in Section~\ref{THF}, with five EPS problems solved but not
% counted because the claim of satisfiability came only after the time limit.
% 
% Most of the systems did better in the EPS problem category.
% While there was no requirement for acceptable solution output in the EPR division, most of the
% systems did output models.
% The question of what constitutes an acceptable model has been a topic of quite hot debate after
% models were deemed unacceptable in the TFN division of \LAST{} (see Section~\ref{TFN}).
% In particular, there are conflicting views of whether or not a saturation is acceptable, e.g.,
% \cite{Pel03-JSC} claims that saturations are: ``useless from a practical point of view''.
% Given that a saturation represents a set of Herbrand models, a salient question is whether that
% set should be a subset of the models of the input formulae.
% Vampire~5.0's preprocessing steps can lead to a saturation whose models is not a subset, but 
% a list of ``Definitions and Model Updates'' is provided to convert the saturation into one 
% whose models are a subset.
% The competition panel ruled that this combination of ``saturation+updates'' constitutes an 
% acceptable model.
% Other types of model output were Vampire's finite models with a domain and mappings, and 
% iProver's Herbrand models expressed as a term algebra.
% The distribution of the different types of models produced is interesting:
% Vampire produced 35 saturations without updates, 26 saturations with updates, and 29 finite 
% models; iProver produced 26 saturations and 57 Herbrand models;
% the other systems produced only saturations.
% Vampire's ability to find finite models gave it a clear advantage.
% 
% As noted in Section~\ref{TPTPProblems}, 191 problems with rating 0.00, 2 problems with rating 
% 0.14, and 11 problems with rating 1.00 were made eligible for the EPR division, all in the EPS
% problem category.
% Of those, 72 problems of rating 0.00, both problems of rating 0.14, and 9 problems of rating 
% 1.00 were selected.
% Of the 72 problems with rating 0.00, 42 were solved by all the systems, 62 were solved by 
% Vampire~5.0, and all the problems were solved by at least one system.
% Of the 9 problems with rating 1.00, only one was solved, by Vampire.
% 
% The individual problem results show that
% 14 problems were unsolved,
% 55 problems were solved by all the systems, and
% 13 problems were solved by only one system, all by Vampire~5.0.
% % Of the 14 unique solutions,
% % 11 were by Vampire,
% % and
% % 1 was by iProver.
% A portfolio approach cannot improve on the individual systems' results.
% 
%--------------------------------------------------------------------------------------------------
\subsection{The UEQ Division}
\label{UEQ}

Table~\ref{UEQSummary} summarizes the results of the UEQ division.
Vampire~5.0.1 won, and as in the FOF division showed significant improvement over Vampire~5.0,
for the same reasons.
The goal directed technique pioneered by Twee~\cite{Sma21} is one of the techniques that led
to Vampire~4.9 winning the UEQ division of CASC-J12 in 2024, and remains in Vampire~5.0.1.
It's value is well appreciated by the Vampire developers, who said: ``tgt-ground {\em [the goal 
directed technique]} is definitely worth having. We are actually considering baking it into the 
default strategy.''
E also tried to adopt the idea, but sadly, as the developer explained: ``my hack for CASC was 
in vain ;-). I carefully added the goal-directed transformation strategy to every search schedule,
but \ldots the preprocessing schedules are now separate data structures, and the transformation 
was only relevant in preprocessing''.

FindProof, mrs, and Prover9 failed to output acceptable proofs for all problems solved.
FindProof simply did not output any proof for 18 problems, and mrs output one proof with syntax 
errors. 
For Prover9, as the developers put it when fixing the output for their next release: ``The bug 
is an amusing one. Prover9 still registers Otter’s legacy infix v (disjunction) as a parse rule, 
so when a TPTP problem legitimately uses \smalltt{v} as a binary functor, the printer rendered 
the legacy in infix notation''.

\begin{table}[tbh]
\caption{UEQ division results}
\label{UEQSummary}
\begin{tabular}{lrrrrrr@{}l@{}rrr}
\hline
ATP System             & Solns & Solved & Avg  & Sot & WC       & \multicolumn{2}{r}{Core} &Uniq & New & Gr'ty \\
                       & /400  & /400   & WC   & AC  & $\mu$Eff & \multicolumn{2}{r}{usage}& /30 & /21 &       \\
\hline
Vampire~5.0.1          &  \bb{370}  &  \bb{370}   & \bb{11.8} & \bb{0.40} & 222     & 4.1       & {\tiny /351} &  \bb{17} &  \bb{12} &  7.9  \\
Twee~2.7               &  286  &  286   & 22.4 & 0.24 & 225     & 6.4       & {\tiny /273} &   2 &   7 &  2.0  \\
CSI++~1.0              &  258  &  258   & 21.5 & 0.19 & \bb{261}     & 5.0       & {\tiny /223} &   1 &   9 & 10.4  \\
E~3.5.1                &  255  &  255   & 20.4 & 0.19 & 244     & 5.3       & {\tiny /226} &   1 &   8 & 12.3  \\
FindProof~0.1          &  251  &  269   & 25.3 & 0.21 & 118     & 6.5       & {\tiny /269} &   1 &   4 &  2.3  \\
Prover9~2026-6A        &  207  &  208   & 21.4 & 0.15 & 202     & 6.2       & {\tiny /179} &   4 &  10 &  \bb{2.0}  \\
iProver~3.9.4          &  198  &  198   & 37.9 & 0.12 &  73     & 5.1       & {\tiny /189} &   1 &   7 & 10.5  \\
Drodi~4.1.1            &  163  &  163   & 37.2 & 0.09 & 127     & \bb{6.7}       & {\tiny \bb{/143}} &   0 &   3 &  4.2  \\
mrs~0.2.0              &   75  &   76   & 36.8 & 0.02 &  33     & 5.8       & {\tiny /71 } &   0 &   2 &  3.8  \\
\hline
\multicolumn{11}{c}{Demonstration division} \\
\hline
\textit{Vampire~5.0}   &  322  &  322   & 16.6 &    - & 295     & 5.2       & {\tiny /310} &   3 &  13 &  8.7  \\
\hline
\end{tabular}
\end{table}

The average times taken and SotAC values are roughly aligned with the division ranking.
CSI++ had a high efficiency because it solved many problems quickly, typically through its
use of E.
As in the FOF and FNT divisions, Drodi and FindProof had high core usage.
Vampire~5.0.1 and Prover9 performed best on the new problems, with Prover9 showing the least
bias towards old problems.
The granularity values are reasonably low, with Prover9 having an optimal value.

The individual problem results (without considering proof output) show that
3 problems were unsolved,
43 problems were solved by all the systems, and
30 problems were solved by only one system.
% (4 problems were solved by only the two versions
% of Vampire, and are counted as unique solutions for Vampire~5.0).
Of the 30 unique solutions,
17 were by Vampire~5.0.1,
6 by Prover9,
3 by Vampire~5.0,
2 by Twee,
and
1 each by CSI++, E, FindProof, and iProver.
A portfolio does not help much, the best being a combination of Vampire~5.0.1 and CSI++~1.0
with 90~s allocated to each, which would solve 373 problems.

%--------------------------------------------------------------------------------------------------
% \subsection{The SLH Division}
% \label{SLH}
% 
% Table~\ref{SLHSummary} summarizes the results of the SLH division.
% Vampire's improved performance on higher-order problems (see Section~\ref{THF}) showed again here.
% As noted in Section~\ref{SLHProblems}, 383 of the 1000 problems had not been solved in the
% testing done before CASC-29, and if many of them were solved in \THIS{} that would indicate
% progress in the field.
% Of the 383 problems, 5 were solved in \THIS{}, 5 by Vampire~5.0, 3 by Vampire~4.9, and 1 by E.
% The lack of progress might be attributed to the lack of tuning that was expected.
% The \THIS{} design informed the developers that the SLH problems would be taken from the same 
% problem set as for CASC-29 and \LAST{}, so that the systems could be tuned using the problems 
% selected for CASC-29 and \LAST{}. 
% According to the systems' developers, no such tuning was done.
% % NOT VAMPIRE ... No, we didn't. The schedule has been fixed for the last two years. It has
% % been trained on the dataset from https://matryoshka-project.github.io/pubs/seventeen.pdf
% % NOT E ... indeed, I did not.
% % NOT LEO-III ... Leo-III is not tuned to anything as a general rule
% 
% The logs from Leo-III revealed that it suffered from the short 15~s CPU time limit.
% The developer explained: ``Leo-III runs in a JVM. 
% The start-up of a JVM uses quite a lot of time, and execution of the Leo-III code is slow at 
% first as the JVM is loading the required classes on-demand. 
% This gets better once proof search is running.''
% A quick experiment was done using a native build that compiles to LLVM code, on two relatively 
% easy problems on which Leo-III timed out in the competition: {\tt SLH0003\verb|^|1} and 
% {\tt SLH0004\verb|^|1}.
% They were solved in 2.4~s and 1.2~s respectively.
% As the developer said: ``Seems to make a difference ... good spot''.
% 
% \begin{table}[tbh]
% \caption{SLH division results}
% \label{SLHSummary}
% \begin{tabular}{lrrrrr}
% \hline
% ATP System             & Solns & Solved & Avg  & Sot  & CPU      \\
%                        & /1000 & /1000  & CPU  & AC   & $\mu$Eff \\
% \hline
% Vampire~5.0            &  433 &   434 &  1.9& 0.16& 317     \\
% E 3.3.0                &  427  &   429  &  2.7 & 0.16& 238      \\
% \hline
% \multicolumn{6}{c}{Demonstration division} \\
% \hline
% \textit{Vampire 4.9}   &    -  &   441  &  2.0 &    - & 311      \\
% cvc5 1.3.0             &    -  &   294  &  3.0 & 0.09 & 117      \\
% Leo-III 1.7.19         &    -  &    38  &  3.3 & 0.01 &  11      \\
% \hline
% \end{tabular}
% \end{table}
% 
% The individual problem results (without considering proof output) show that
% 479 problems were unsolved,
% 21 problems were solved by all the systems, and
% 102 problems were solved by only one system (41 problems were solved by only the two versions
% of Vampire, and are counted as unique solutions for Vampire~5.0).
% In \LAST{} 258 problems were solved by all the systems: Vampire~4.9, E~3.2.0, cvc5~1.1.3, and
% the CASC-29 winner E~3.1.
% The impact of Leo-III in \THIS{} is evident -- 234 problems were solved by all the other systems.
% Of the 102 unique solutions,
% 45 were by Vampire~5.0,
% 36 by E,
% 13 by cvc5,
% 7 by Vampire~4.9,
% and
% 1 by Leo-III
% A portfolio approach works well here --
% a portfolio of E, Vampire~4.9, and cvc5, with 5~s allocated to each, would solve 482 problems.
% 
% The SLH division has served it's purpose, to raise awareness of the needs of the Sledgehammer
% system -- higher-order problems to be solved within a short CPU time limit.
% Sadly it has not inspired developers to tune to the problem set.
% The SLH division will go on hiatus from \NEXT{}.

%--------------------------------------------------------------------------------------------------
% \subsection{The ICU Division}
% \label{ICU}
% 
% Table~\ref{ICUSummary} summarizes the results of the ICU division, and Table~\ref{ICUMatrix}
% shows how many of the problems submitted by each system's developer were solved by each
% system (the columns correspond to the problem submissions from the corresponding system row).
% As might have been expected, the Vampires were dominant (see Section~\ref{FOF} for the
% reasons).
% Most of the systems solved all or most of their submitted problems, and Vampire~5.0 additionally 
% solved all of Drodi's problems.
% CSI\_Enigma and E also solved a reasonable cross-section of the other systems' problems.
% 
% \begin{table}[tbh]
% \caption{ICU division results}
% \label{ICUSummary}
% \begin{tabular}{lrrrrrr@{}l@{}}
% \hline
% ATP System             & Solns & Solved & Avg  & Sot  & WC       & \multicolumn{2}{r}{Core}  \\
%                        & /100  &  /100  & WC   & AC   & $\mu$Eff & \multicolumn{2}{r}{Usage} \\
% \hline
% Vampire~5.0            &   69 &    69 & 58.0 & 0.39& 172     & 5.9       & {\tiny /66}  \\
% E 3.3.0                &   42  &    42  &115.0 & 0.19 &  43      & 6.7       & {\tiny /41}  \\
% iProver 3.9.3          &   28  &    34  &114.1 & 0.14 &  17      & 7.1       & {\tiny /34}  \\
% Drodi 4.1.0            &   16  &    21  & 50.5& 0.07 &  41      & 7.2      & {\tiny /20}  \\
% CSE\_E 1.7             &   16  &    18  &143.5 & 0.06 &   8      & 2.8       & {\tiny /18}  \\
% Connect++ 0.6.1        &    0  &     1  &417.5 & 0.00 &   0      & 0.0       & {\tiny /0}   \\
% \hline
% \multicolumn{6}{c}{Demonstration division} \\
% \hline
% \textit{Vampire~4.9}   &    -  &    70  & 21.7 &    - & 195      & 5.9       & {\tiny /67}  \\
% CSI\_Enigma 1.0.6      &    -  &    46  &124.2 & 0.21 &  28      & 6.3       & {\tiny /46}  \\
% \hline
% \end{tabular}
% \end{table}
% 
% \begin{table}[tbh]
% \caption{ICU matrix}
% \label{ICUMatrix}
% \begin{tabular}{lrrrrrrrr}
% \hline
% ATP System             & VAM & EEE & IPR & DRO & CSE & CPP & VVA & CSI \\
% ATP System             & /13 & /13 & /13 & /13 & /13 & /12 & /10 & /13 \\
% \hline
% Vampire~5.0            &  13&   7 &  11 &  13&   7 &   0 &   6&  11\\
% E 3.3.0                &   1 &  13&   3 &   6 &   9&   0 &   0 &  10 \\
% iProver 3.9.3          &   2 &   2 &  12&  10 &   1 &   0 &   0 &   7 \\
% Drodi 4.1.0            &   0 &   2 &   3 &  12 &   1 &   0 &   0 &   3 \\
% CSE\_E 1.7             &   0 &   4 &   4 &   1 &   3 &   0 &   0 &   5 \\
% Connect++ 0.6.1        &   0 &   0 &   0 &   1 &   0 &   0 &   0 &   0 \\
% \hline
% \multicolumn{9}{c}{Demonstration division} \\
% \hline
% \textit{Vampire~4.9}   &   8 &   6 &  12 &  13 &  11 &   0 &  10 &   9 \\
% CSI\_Enigma 1.0.6      &   1 &   9 &   6 &   9 &  11 &   0 &   1 &   9 \\
% \hline
% \end{tabular}
% \end{table}
% 
% The SotAC values align with the division ranking, but the efficiency values are less aligned.
% iProver and CSI\_Enigma had lower efficiencies due to their higher average times taken.
% Drodi's performance stands out thanks to its low average time taken and high core usage.
% 
% The individual problem results (without considering proof output) show that
% 14 problems were unsolved,
% no problems were solved by all the systems, and
% 24 problems were solved by only one system (14 problems were solved by only the two versions
% of Vampire, and are counted as unique solutions for Vampire~5.0).
% Of the 24 unique solutions,
% 18 were by Vampire~5.0,
% 5 by Vampire~4.9,
% and
% 1 by E.
% A portfolio approach naturally works well here because the problems were purposefully chosen
% for each system.
% A combination of the two Vampires and E, with 160~s allocated to each, would solve 86 problems.

%--------------------------------------------------------------------------------------------------
% \subsection{The LTB Division}
% \label{LTB}
%
% Table~\ref{LTBSummary} summarizes the results of the LTB division.
% The Vampires solved the most problems, demonstrating the efficacy of their dynamic strategy
% tuning:
% The idea is to start by attempting all versions of all problems with a short time limit, to get
% an initial set of problems solved.
% Following that an auto-tuning mode starts, repeatedly iterating through the remaining unsolved
% problems, randomly selecting a variant of an unsolved problem and randomly selecting a strategy
% for the version, weighted according to how successful the strategies have been so far.
% The strategy weighting is updated dynamically by giving a successful strategy a weight
% proportional to how many previous strategies failed to solve the problem version, and giving
% an unsuccessful strategy a constant penalty.
% Over time the successful strategies are given higher priority, and different (successful)
% strategies are used in subsequent attempts on unsolved problems.
% For \THIS{} Vampire~4.7 added a slight variation that introduced a fresh strategy with a
% 1\% probability at each attempt, by randomly permuting an existing strategy.
% If the fresh strategy solved the problem it was kept in the pool of strategies, otherwise it
% was discarded.
% The purpose was to inject some randomness into very long runs, but the developer reported that
% he has ``no idea if it actually helped!''.
%
% The selection of problems and problem versions was evidently harder for the systems than in
% \LAST{}, where Vampire~4.6 solved 78\%, iProver~3.5 70\%, and E~2.6 68\% of the problems,
% compared to 47\%, 43\%, and 41\% respectively, this year.
% Polymorphism continues to be a challenge for the ATP systems -- see the \LAST{} report
%~\cite{SD22-CASC} and the discussion of Table~\ref{VersionCounts} below.
% % Recall from Section~\ref{LTBProblems} that there were 490 problems that had not been proved
% % during the problem testing before the competition.
% In \LAST{} Vampire~4.6 solved 262 problems by finding that they had contradictory axioms.
% This year Vampire~4.7 solved 146 by contradictory axioms.
% Of the 146, 41 were TF0, 30 TF1, 66 TX0, 5 TH0, and 4 TH1.
% Contradictory axioms is a recurring ``feature'' of problems generated by Sledgehammer, where
% contradictory axioms can appear in proofs of negation, proofs by contradiction, and proofs by
% case analysis where some cases do not apply.
%
% \begin{table}[htb]
% \caption{LTB division results}
% \label{LTBSummary}
% \begin{tabular}{lrrrrr}
% \hline
% ATP System           & LTB  &  Avg & Sot  &  WC &Core \\
%                      &/8000 &  WC  & AC   &$\mu$Eff&usage \\
% \hline
% Vampire~4.7          & 3729 & 46.3 & 0.09 & 154 & 5.7  \\
% iProver 3.6          & 3446 & 50.1 & 0.05 & 126 & 6.1  \\
% E 3.0                & 3292 & 52.5 & 0.05 & 116 & 1.6  \\
% \hline
% \multicolumn{6}{c}{Demonstration division} \\
% \hline
% \textit{Vampire~4.6} & 3725 & 46.4 &    - & 163 & 5.7  \\
% \hline
% \end{tabular}
% \end{table}
%
% Table~\ref{VersionCounts} shows the numbers of each version of the problems solved.
% There were 4532 problems solved across all systems and problem versions.
% There is quite a high degree of specialization to specific problem versions, with Vampire~4.7
% having the broadest range.
% The most problems were solved in the lowest logic, with iProver solving all its 3446 problems
% in TF0.
% The benefit of being able to solve problems in higher logics is clear, and there appears to be
% some space for improvement by adding TX1 capability.
%
% \begin{table*}[hbt]
% \caption{Problems solved for each version}
% \label{VersionCounts}
% \begin{tabular}{c|R{8mm}R{8mm}R{8mm}R{8mm}|r}
% \hline
% Ver    & V~4.7& V~4.6& iPro & E    & Union \\
% \hline
% TF0    &  959 & 1948 & 3446 & 2816 &  4008 \\
% TF1    &  942 & 1307 &    0 &    0 &  1870 \\
% TX0    & 1541 &    0 &    0 &    4 &  1544 \\
% TX1    &    0 &    0 &    0 &    0 &     0 \\
% TH0    &  260 &  449 &    0 &  472 &   999 \\
% TH1    &   27 &   21 &    0 &    0 &    41 \\
% \hline
% Sum    & 3729 & 3725 & 3446 & 3292 &  4532  \\
% \hline
% \end{tabular}
% \end{table*}
%
% The Vampires' dynamic strategy tuning approach meant that most problems and problem versions
% were attempted more than once; only 40 problems were attempted (and thus solved) only once, six
% problems were solved in the 77th attempt, and one problem was attempted (but not solved) 78 times.
% While over half of the solutions were found by 12 attempts, the figure shows how the number of
% problems solved keeps increasing, up to the six solved in the 77th attempt.
% The multiple attempts on problems is also reflected in Vampire~4.7's accumulated CPU times across
% the solved problems, from 1155 problems accumulating less than 1~s of CPU time, up to 2 problems
% that accumulated just over 4000~s of CPU time -- see the performance graph for the LTB
% division\footnote{%
% \url{https://tptp.org/CASC/J11/WWWFiles/ResultsPlots.html\#LTBProblems}}.
% Vampire~4.6 took the same approach as Vampire~4.7 but without the new strategy permutation
% feature.
% Vampire~4.6 was disadvantaged by the different versions of problems available in \THIS{}:
% In \LAST{} there were first-order form (FOF) problems, and no TXF problems.
% As a result this year Vampire~4.6 spent time looking for FOF problems that did not exist,
% ignored the TXF problems because TXF was not part of \LAST{}, and attempted the THF problems to
% a greater extent than Vampire~4.7, as is evident in Table~\ref{VersionCounts}.
% E also attempted more than one version of each problem, by first trying the TF0 versions, then
% the TX0 versions of the remaining unsolved problems, and finally the TH0 versions of the
% remaining still unsolved problems.
% E used a dynamic scheduling algorithm that gave each attempt the remaining time divided by the
% number of remaining versions of unsolved problems.
% E solved only a very small number of TX0 problems; the developer believes this was caused by E
% converting TX0 problems to TF0 problems in preprocessing, resulting in problems that were
% nearly identical to the TF0 problems on which the strategies had already failed in the last
% iteration.
% iProver attempted only the TF0 versions of the problems.
%
% % \begin{figure}[ht!]
% % \begin{center}
% % \includegraphics[width=0.4\columnwidth]{V47AttemptsSolvedPlot}
% % \caption{Vampire~4.7 problems solved by attempts}
% % \label{V47AttemptsSolvedPlot}
% % \end{center}
% % \end{figure}
%
% The individual problem results show that
% 3468 problems were unsolved,
% 2552 problems were solved by all the systems,
% and 919 problems were solved by only one system
% (263 problems were solved by only the two versions of Vampire, and are counted as
% unique solutions for Vampire~4.7).
% Of the 919 unique solutions,
% 425 were by Vampire~4.7,
% 175 were by Vampire~4.6,
% 169 by iProver,
% and
% 150 by E.
% As the systems have different approaches to solving the problems, a portfolio approach that shares
% the available time cannot be (easily) considered.
% However, Table~\ref{VersionCounts} shows that giving all systems the full 172800~s solves 4532
% problems.
% Interestingly, just Vampire~4.6 and iProver~3.6 together would solve 4141 problems, and adding
% Vampire~4.7 to that mix would solve 4382 problems.
%
%--------------------------------------------------------------------------------------------------
\section{System Descriptions}
\label{Systems}

Each year the competition report features short system descriptions of systems that stand out in
some way.
For \THIS{} the salient systems are the newcomers that performed relatively well (as newcomers)
in the FOF division: FindProof~0.1, SUPr~1.0, and mrs~0.2.0.
These short descriptions of the systems were written by their entrants.

\textbf{FindProof} is 

\textbf{SUPr}~\cite{KP26} is a saturation-based theorem prover built as a native component of the 
Sigma knowledge engineering environment (Rust implementation -- SigmaKEE-rs)~\cite{PS14}.
SUPr is optimized to reason directly over the Suggested Upper Merged Ontology (SUMO)~\cite{Pea11}, 
with additional support to solve standalone TPTP problems. 
The core search procedure is a given-clause saturation loop implementing ordered resolution and 
superposition with selective literal selection, forward and backward demodulation, and a 
Knuth-Bendix reduction ordering. 
Clauses and their constituent terms are represented in a content-addressed, hash-consed arena. 
This content based addressing scheme is structured to allow batched, hardware accelerated 
unification via binary arithmetic.
As SUPr is built specifically for reasoning over SUMO, SUPr's proving strategy depends on heavy 
caching of supporting evidence for new axioms at the time of ingest, and use-case-specific 
semantic decoding.
Novel pre-caching of SUMO Inference Engine~\cite{HV11} (SInE) occurences and trigger indices allows 
for fast per-problem axiom selection. 
Axioms are additionally ranked by structural relevance to the conjecture~\cite{LX22}, and those 
that rank highly but were dropped by SInE are rescued and reintroduced to the problem.
Search-loop tuning is organized as a small number of strategies running as concurrent threads 
that race for a solution. 
Threads share a single clause stack that remains unconflicted due to the atomicity of the 
content-based hashing scheme.
 
\textbf{mrs}~\cite{Rol26} is is an automated theorem prover for first-order logic with equality. 
It implements the superposition calculus~\cite{BG94} in a given-clause loop, with ordered 
resolution, factoring, equality resolution, equality factoring, and paramodulation, all oriented 
by KBO or LPO with rarity-based symbol precedence. 
Non-Horn clauses are split via AVATAR~\cite{Vor14} using the CaDiCaL~\cite{BF+24} SAT solver.
mrs is written in Rust.
The search runs a parallel portfolio, with each strategy maintaining its own clause set. 
A central design feature is proof-preserving cross-strategy sharing: workers exchange selected 
positive unit equalities with complete ancestor chains, remap those chains into the receiving proof 
store, and use only the equality heads in subsequent search.
Rust's ownership and thread-safety checks make the synchronized shared pool explicit and race-safe.
The portfolio for each category division (FNE, FEQ, UEQ) was selected by running all 15 base 
strategies solo and using greedy set-cover analysis to choose an eight-strategy portfolio. 
Local analysis identified clause generation and redundancy, rather than raw iteration speed, as 
the main bottleneck. 
The competition's objective scores and rapid feedback loop made profiling, debugging, and heuristic 
experiments focused, repeatable, and fun, directly shaping subsequent work on inference 
restrictions, resource controls, and proof-preserving sharing. 
CASC pushed mrs to fix real bugs and close gaps that show up only at scale.

%--------------------------------------------------------------------------------------------------
\section{Acceptable Proofs}
\label{Acceptable}

In Sections~\ref{THF}, \ref{FOF}, and~\ref{UEQ}, the extent to which the systems did not
output acceptable proofs is discussed.
For newcomers to CASC this is quite understandable, as generating proofs is a non-trivial task.
For the old timers, the increasing emphasis on proof production since CASC-J12 should be
having a positive effect.
This section examines the requirements for acceptable proofs for the last three CASCs -- CASC-J12,
\LAST{}, and \THIS{}, and evaluates the progress made between \LAST{} and \THIS{} for systems
that were entered in both (with an updated version in \THIS{}).

%--------------------------------------------------------------------------------------------------
\subsection{Requirements}
\label{Requirements}

Prior to \LAST{} the CASC panel was the only adjudication of whether systems' proofs were
acceptable. 
The following guidelines were set for CASC-J12~\cite{Sut24-CASC-J12}:
\begin{packed_enumerate}
\item[1.] Proofs must be complete, starting at formulae from the problem, and ending at the 
      conjecture (for axiomatic proofs) or a false formula (for proofs by contradiction, 
      e.g.,~CNF refutations).
\item[2.] For proofs that use translations from one form to another, e.g.,~translation of FOF 
      problems to CNF, the translations must be adequately documented.
\item[3.] Proofs must show only relevant inference steps.
\item[4.] Inference steps must document the parent formulae, the inference rule used, and the 
      inferred formula.
\item[5.] Inference steps must be reasonably fine-grained.
\end{packed_enumerate}
The CASC-J12 report~\cite{Sut25-CASC} noted the benefits of moving towards a ``verified track'' 
in CASC.
An incremental approach was adopted, and the following requirements for acceptability were added 
in \LAST{}:
\begin{packed_enumerate}
\item[6.] Proofs must be in the TPTP format.
\item[7.] Annotated formulae in a proof must be uniquely named.
\item[8.] Proofs must be structurally correct, e.g., derivations must be acyclic.
\item[9.] Proofs that negate the \smalltt{conjecture} must give the resultant formula the role 
      \smalltt{negated\_conjecture}, and annotate the step as \smalltt{status(cth)}.
\item[10.] Proofs that introduce new symbols in definitions must provide an \smalltt{introduced()} 
      record to document the new symbol.
\end{packed_enumerate}
The TPTP4X parser~\cite{Sut10} and GDV proof verifier~\cite{Sut06} were made available to the 
systems' developers so that they could check their proofs.
Note that there was no requirement for logical correctness, and no proof verification was done.

The impact was interesting, and the data showed that the ATP systems still needed to improve 
their proofs.
A particular phenomenon that was observed was that many of the ATP systems' proofs included
``inference steps'' that were calls to SAT/ATP/SMT systems, which typically inferred {\tt \$false}
to end the ATP system's refutation.
These final calls to external ATP systems were typically reported with all the clauses sent to 
the external system listed as parents -- the most extreme case listed 72367 parent clauses that 
were sent to the external Z3 system.
A limited analysis revealed that typically only a subset of the listed parents were used by the 
SAT/ATP/SMT system to infer {\tt \$false}.

To encourage developers to improve their systems' proofs, and avoid the second weakness, the 
requirements for acceptable proofs were refined in \THIS{}:
\begin{packed_enumerate}
\item[11.] Inference steps must list exactly those parents that are used in the inference.
      The granularity measure was added to the results to track the maximal numbers of 
           parents in the inference steps of proofs.
\item[12.] Proofs that make \smalltt{assumption}s must propagate and eventually discharge the 
      assumptions.
\end{packed_enumerate}

%--------------------------------------------------------------------------------------------------
\subsection{Progress}
\label{Progress}

\begin{table}[htb]
\caption{Progress in acceptable proof output}
\label{ProofProgress}
\begin{tabular}{l|lrrr|lrrr|rr}
\hline
            &\multicolumn{4}{l|}{\LAST{}}         & \multicolumn{4}{l|}{\THIS{}}       & $\Delta$ Loss &Progress \\
System      & Version  & Solved & Solns &  Loss   & Version & Solved & Solns &  Loss  &               &         \\
\hline
            &\multicolumn{4}{l|}{THF/500}         & \multicolumn{4}{l|}{THF/400}       &               &          \\
E           & 3.3.0    &   395 &    375 &  5.06\% & 3.5.1   &   324 &    323 &  0.31\% &       -4.75\% &   3.75\% \\
Vampire     & 5.0      &   449 &    439 &  2.23\% & 5.0.1   &   352 &    340 &  3.41\% &        1.18\% &  -1.00\% \\
Leo-III     & 1.7.19   &   242 &    241 &  0.41\% & 1.8.0   &   160 &    137 & 14.38\% &       13.96\% &  -5.55\% \\
\hline
            &\multicolumn{4}{l|}{FOF/500}         & \multicolumn{4}{l|}{FOF/400}       &               &          \\
Drodi       & 4.1.0    &   325 &    210 & 35.38\% & 4.1.1   &   215 &    214 &  0.47\% &      -34.92\% &  22.75\% \\
iProver     & 3.9.3    &   367 &    278 & 24.25\% & 3.9.3   &   250 &    222 & 11.20\% &      -13.05\% &  10.80\% \\
Connect++   & 0.6.1    &   102 &      0 &100.00\% & 0.7.2   &    64 &     24 & 62.50\% &      -37.50\% &  10.40\% \\
E           & 3.3.0    &   364 &    364 &  0.00\% & 3.5.1   &   264 &    264 &  0.00\% &             - &   6.80\% \\
SPASS-SCL   & 0.1      &    11 &      0 &100.00\% & 0.1.1   &    14 &     14 &  0.00\% &     -100.00\% &   2.20\% \\
Vampire     & 5.0      &   455 &    455 &  0.00\% & 5.0.1   &   381 &    380 &  0.26\% &        0.26\% &  -0.25\% \\
\hline
            &\multicolumn{4}{l|}{UEQ/300}         & \multicolumn{4}{l|}{UEQ/400}       &               &          \\
iProver     & 3.9.3    &   205 &     62 & 69.76\% & 3.9.3   &   198 &    198 &  0.00\% &      -69.76\% &  47.67\% \\
Twee        & 2.6.0    &   258 &    258 &  0.00\% & 2.7     &   286 &    286 &  0.00\% &             - &  14.50\% \\
E           & 3.3.0    &   222 &    222 &  0.00\% & 3.5.1   &   255 &    255 &  0.00\% &             - &  10.25\% \\
Vampire     & 5.0      &   263 &    263 &  0.00\% & 5.0.1   &   370 &    370 &  0.00\% &             - &  -4.83\% \\
Drodi       & 4.1.0    &   121 &    121 &  0.00\% & 4.1.1   &   163 &    163 &  0.00\% &             - &  -0.42\% \\
\hline
\end{tabular}
\end{table}

\[
Loss\,\text{\em \%} = 
\dfrac{Solved - Solns}{Solved}
\]

\[
Progress = 
\begin{cases}
\text{if 0\% Loss both years} , & \text{\em \LAST{}}\!\left(\dfrac{Solved}{Problems}\right) - \text{\em \THIS{}}\!\left(\dfrac{Solved}{Problems}\right) \\[10pt]
\text{otherwise}              , & \text{\em \LAST{}}\!\left(\dfrac{Solved - Solns}{Problems}\right) - \text{\em \THIS{}}\!\left(\dfrac{Solved - Solns}{Problems}\right)
\end{cases}
\]

There were two particular points of weakness that were observed in the \LAST{} proofs.
First, 459 of the 3961 proofs, i.e.,~11.5\%, from the THF, FOF, and UEQ divisions had syntax or 
structural flaws.
The situation improved in \THIS{} with only 200 of 5065 proofs, i.e., 3.9\%, having those types
of flaws.

After \THIS{} some further detailed analysis of the proofs revealed that many had ill-formed
documentation of inference steps, in the \smalltt{introduced{}} and \smalltt{inference()} records.
The \THIS{} panel ruled that the documentation of the required format for these records was 
inadequate, which allowed those proofs to be counted.
The documentation has been improved for these two points and more, and more precise requirements 
have became possible.

The granularity values reported in Sections~\ref{THF}, \ref{FOF}, \ref{FNT}, and~\ref{UEQ} provide
some indication of the situation in \THIS{}.
As a counter point to the maximal value of 72367 in \LAST{}, the maximal value in \THIS{} was
12651.

%--------------------------------------------------------------------------------------------------
\subsection{The Future}
\label{FutureProofs}

Discussions during the IJCAR-13 conference that hosted \THIS{} led to the conclusion that 
requirement 11 is unrealistic because precisely extracting the set of parents that were used 
by a SAT/ATP/SMT system might not be possible in all cases.
From the organizational side, checking conformance to the requirement might not be possible in 
all cases. 
It became clear that the systems can only ``aim to'' comply, leading to a change to the 
requirement, specified in requirement 11' below.
Even with a minimization of the parents list in inference steps, it is debatable whether a 
successful call to a SAT/ATP/SMT system is really an ``inference step''.
The SAT/ATP/SMT systems typically build a proof of their own, and that proof structure is not 
shown in the ATP system's proof that reports the step.
It would be preferable for the SAT/ATP/SMT system's proof to be exposed, and grafted into the
ATP system's proof.
Grafting in the proofs of subsystems will be encouraged in \NEXT{}, specified in requirement
14 below.

The requirements for \NEXT{} will include:
\begin{packed_enumerate}
\item[9'.] Proofs that negate the conjecture must correctly annotate the step as
      \smalltt{status(cth)} and have a single parent with the role \smalltt{conjecture}.
\item[10'.] New symbols introduced in derivations must be documented in a \smalltt{new\_symbols()}
      record. 
      The new symbols listed in a \smalltt{new\_symbols()} record must be unique and globally new.
\item[11'.] Inference steps must aim to list exactly those parents that are used in the inference. 
      Inference steps may list parents that are not necessary for the inference\footnote{%
      For example, in the inference of \smalltt{p(a)} from the parents \smalltt{p(a)} and 
      \smalltt{a = a}, the equality is not necessary but might be used in an (unnecessary) 
      equality substitution step, but in contrast \smalltt{p(b)} would probably not be used 
      as a parent in that inference step.}.
\item[13.] Proofs should not use a TPTP language that is more expressive than that of the problem.
\item[14.] Proofs should graft in subproofs from external systems.
\item[15.] Particular requirements for symbol definitions, introduced tautologies, theory axioms, 
      and Skolemization steps will be provided.
\end{packed_enumerate}
The BNF-based parser, the TPTP4X parser, and an improved version of GDV will be made available 
to the systems' developers so that they can check their proofs.
These more rigorous requirements for proofs in CASC resonate with the ProoVer proof verifier 
competition that was held as a companion competition to \THIS{}.
ProoVer is described in Section~\ref{ProoVer}.

Syntax checking of models will be added in \NEXT{}, as explained in Section~\ref{Acceptable}.

MAYBE USE THIS. One thing which you could say about the LisaST if you have space and find it interesting, is that it is, if I'm not mistaken, the first solver to be reasonably competitive at CASC that can produce proofs that directly checkable by a small kernel/proof assistant, including clausification :)

%--------------------------------------------------------------------------------------------------
\subsection{What is a Proof?}
\label{WhatIsAProof}

The question ``what is a proof'' has been discussed in great detail from many 
perspectives~\cite{JEFF}.
In ATP a common proof output is a refutation, which might in itself not be considered a proof
{\em of} a theorem (a formula that was a conjecture), but rather a proof that a conjecture 
formula {\em is} a theorem.
Consider these contrasting examples.
First a proof of a theorem:
\begin{quote}
Rule: If it is raining, take your umbrella.\\
Fact: It's raining.\\
Conjecture: Take your umbrella.\\
Proof: The fact makes the rule apply, so accept the conclusion of the rule.\\
Theorem: Take your umbrella.
\end{quote}
Now a proof that the conjecture is a theorem:
\begin{quote}
Rule: If it is raining, take your umbrella.\\
Fact: It's raining.\\
Conjecture: Take your umbrella.\\
Proof: Imagine you don't take your umbrella.
Then from the contrapositive of the rule, it's not raining.
That contradicts the fact that it's raining.
Thus it is unimaginable that you don't take your umbrella.\\
Theorem: Take your umbrella.
\end{quote}

Should I define my general axiomatic proof notion here?
In the long term ATP systems might strive to produce proofs {\em of} theorems, rather than
only proofs that conjectures {\em are} theorems.

%--------------------------------------------------------------------------------------------------
\section{The ProoVer Competition}
\label{ProoVer}

The Proof Verification Competition (ProoVer) is a new competition, evaluating proof
checkers, i.e.,~systems that verify whether a given proof is correct wrt its conjecture. 
ProoVer 2026 was the first edition, held as a companion competition to \THIS{}, sharing its 
computers and much of its organizational machinery.
ProoVer 2026 was organized by Julie Cailler and Simon Guilloud, and overseen by the same 
panel as \THIS{}: Cl\'audia Nalon, Aart Middeldorp, and Martina Seidl.
Full details of the competition design, the entrant systems, and the results are given in the 
ProoVer 2026 proceedings~\cite{ProoVer2026}. 
This section gives an overview of the competition.

A ProoVer entrant is a proof checker that, given an untyped first-order (FOF) TPTP problem and a
corresponding TSTP derivation~\cite{SS+06}, must decide whether the derivation is a valid
refutation of the problem.
A derivation consists of a sequence of syntactically correct TSTP inference steps.
The competition problem set was made of $50$ correct proofs and $50$ incorrect (aka ``evil'') 
ones that had been altered to become incorrect.
Most steps can be discharged to an external ATP system, but the proof checker must verify a small
set of specified rules (Skolemization, the negated conjecture, and correspondence between the 
problem file and the proof file) internally.

Ten proof checkers competed: eight in the ranked regular division, and two in the unranked 
demonstration division (due to the need for internet access).

The ProoVer 2026 problem set comprised $100$ proof-validation problems, each a valid (Good) or 
a flawed (Bad) TSTP derivation. 
However, $9$ problems were later cancelled as ambiguous or non-compliant, leaving 91 (44~Good, 
47~Bad) that counted for the ranking. The possible score ranged therefore from $138$ to $-514$.
The systems are ranked according to the following grading scheme:

\begin{itemize}
    \item Identifying a bad proof as bad = $+2$
    \item Identifying a good proof as good = $+1$
    \item Giving up = $0$
    \item Identifying a good proof as bad = $-1$
    \item Identifying a bad proof as good = $-10$
\end{itemize}


Table~\ref{ProoVerSummary} lists the entrants, grouped by division, together with their final 
score. 
GAPT~2.20 won the regular division and, having a student developer, also won the 
\emph{ProoVer 2026 Best Student Proof Checker} award.
VaLeaDate~0.1, which had the best record at identifying Bad proofs, won the 
\emph{ProoVer 2026 Safest Proof Checker} award.

\begin{table*}[tbh]
\caption{ProoVer 2026 entrants and final ranking}
\label{ProoVerSummary}
\begin{tabular}{llr}
\hline
Proof Checker & Entrant & Score \\
\hline
\multicolumn{3}{c}{\textit{Regular division}} \\
\hline
GAPT~2.20         & Fabian Achammer \& Martin Riener         & 111 \\
VaLeaDate~0.1     & Jonas Bodingbauer \& Laura Kovacs        & 98  \\
N\"orgler~1.1     & Alexander Steen et al.                   & 93  \\
ProofCheck~1.0    & Jeff Machado \& Larry Lesyna             & 75  \\
ProofGuard~1.0    & Matthew Farah \& Alicia Marinache        & 49  \\
mrs-proover~0.2.0 & Olivier Roland                           & 47  \\
PyCheck~0.1       & Stephan Schulz                           & 41  \\
CheckProof~0.1    & Nik Murzin                               & $-103$ \\
\hline
\multicolumn{3}{c}{\textit{Demonstration division}} \\
\hline
GDV~2.0           & Geoff Sutcliffe                         & $-4$ \\
GDV-LP~2.0        & Fr\'ed\'eric Blanqui \& Geoff Sutcliffe & $-6$ \\
\hline
\end{tabular}
\end{table*}

The ProoVer 2026 problem set is available online, together with an explanation of each problem.
ProoVer's first edition led to multiple discussions, serving the broader goal of improving the 
TSTP specification.

%--------------------------------------------------------------------------------------------------
\section{Conclusion}
\label{Conclusion}

\THIS{} was the \NTHCOMPETITION{} annual evaluation of fully automatic, classical logic, Automated
Theorem Proving (ATP) systems.
% The competition provided exposure for ATP systems both within and beyond the
% ATP community, and provided an overview of the implementation state of fully
% automatic classical logic ATP systems.
\THIS{} fulfilled its objectives by evaluating the relative capabilities of ATP systems,
and stimulating development and interest in ATP.
The highlight of \THIS{} was the structural checking of proofs, which revealed flaws in many
systems proofs.
This will lead to improvements in proof output over the next year.
Despite the continued dominance of the Vampire systems, the developer and user communities have
remained engaged with CASC.
A key to this is CASC's stimulating and interactive nature: a live competition run during one day 
of the conference; opportunities for users and observers to be involved; a social dinner where 
entrants and associates can interact; distinctive T-shirts that identify system developers and 
others who are part of the competition.

While the design of CASC is mature and stable, each year's experiences lead to ideas for changes
and improvements.
Changes being planned for \NEXT{} are:
\begin{packed_itemize}
\item Only two installation packages will be allowed for each entry -- one for theorem proving/
      unsatisfiability checking, and one for countermodel finding/satisfiability checking.
      The organizer will provide template \smalltt{starexec\_run} packages to smooth out the
      creation and submission of working installation packages.
\item All solutions (proofs and models) must be in TPTP syntax, error-free according to TPTP4X.
\item Proof checking will be extended to include:
      \begin{packed_itemize}
      \item Leaves must come from the problem.
      \item \smalltt{inference} and \smalltt{introduced} records must be correctly formatted.
      \item All new symbols must be correctly and unique documented in an \smalltt{introduced} 
            record.
      \item Skolemization steps must be correctly documented in an \smalltt{inference} record.
      \end{packed_itemize}
\item Models must be in TPTP syntax.
\item Syntax checking of solutions will be done at runtime during the competition.
\item No sample proofs will be requested, to avoid any mistaken belief that an approved sample
      proof implies that all proofs will be accepted.
\end{packed_itemize}

As always, the ongoing success and utility of CASC depends on ongoing contributions of problems
to the TPTP.
The automated reasoning community is encouraged to continue making contributions of all types of
problems.

%--------------------------------------------------------------------------------------------------
% \begin{acks}
% This material is based upon work supported by the National Science Foundation
% under Grant No. 1730419.
% \end{acks}
%--------------------------------------------------------------------------------------------------
\bibliographystyle{ios1}
\bibliography{Bibliography}
%--------------------------------------------------------------------------------------------------
\end{document}
%--------------------------------------------------------------------------------------------------
